\documentclass[lettersize,journal]{IEEEtran}
\usepackage{algorithmic}
\usepackage{algorithm}
\usepackage{array}
\usepackage[caption=false,font=normalsize,labelfont=rm,textfont=rm]{subfig}
\usepackage{textcomp}
\usepackage{stfloats}
\usepackage{url}
\usepackage{verbatim}
\usepackage{cite}
\usepackage[acronym]{glossaries}
\usepackage{color}
\usepackage{amsthm,amsmath}
\usepackage{graphicx}
\usepackage{amsmath,amssymb,amsfonts}
\usepackage{bm}
\usepackage{mathrsfs}
\usepackage{indentfirst}
\usepackage{multirow}
\usepackage{float}
\usepackage{booktabs}
\usepackage{tabularx}
\usepackage{mhchem}
\usepackage[figuresright]{rotating}
\usepackage[switch]{lineno}
\usepackage{threeparttable, tabularx}
\usepackage{bbding}
\usepackage{makecell}
\usepackage{arydshln}
\usepackage{xcolor}
\usepackage{tikz}
\usepackage[export]{adjustbox}
% \hyphenation{op-tical net-works semi-conduc-tor IEEE-Xplore}
\usepackage{subfiles} % Best loaded last in the preamble
\usepackage{xcolor} % 引入颜色宏包
\hyphenation{op-tical net-works semi-conduc-tor IEEE-Xplore}
% updated with editorial comments 8/9/2021

\renewcommand{\thepage}{}

% 定义一个名为 \rev 的命令
% 使用方法：\rev{这是修改的内容}
\newcommand{\rev}[1]{\textcolor{blue}{#1}}

\begin{document}

\title{\rev{OOD-Robust Degradation Modeling for Vehicular Battery Management Systems} via Uncertainty-Driven Feature Manipulation with Topological Constraint}

\author{
    Xinhui~Lin,
    Yuanguo~Bi,~\IEEEmembership{Member,~IEEE},
    Rao~Fu,
    Qiang~He,
    Jun~Liu,~\IEEEmembership{Member,~IEEE},
    and Liang~Zhao,~\IEEEmembership{Member,~IEEE}
    \vspace{-1.5em}
        % <-this % stops a space
% \thanks{This paper was produced by the IEEE Publication Technology Group. They are in Piscataway, NJ.}% <-this % stops a space
% \thanks{Manuscript received April 19, 2021; revised August 16, 2021.}
\thanks{(\emph{Corresponding author: Yuanguo Bi.})}
\thanks{Xinhui~Lin, Yuanguo~Bi, and Rao~Fu are with the School of Computer Science and Engineering, Northeastern University, Shenyang 110169, China, and also with the Engineering Research Center of Security Technology of Complex Network System, Ministry of Education, Shenyang 110169, China (e-mail: linxinhui@stumail.neu.edu.cn; biyuanguo@mail.neu.edu.cn; 2010764@stu.neu.edu.cn).}
\thanks{Qiang~He is with the School of Computer Science and Engineering, Northeastern University, Shenyang 110169, China (e-mail: heqiang@mail.neu.edu.cn).}
\thanks{Jun~Liu is with the School of Electronic and Information Engineering, Beihang University, Beijing 100191, China (e-mail: liujun2019@buaa.edu.cn).}
\thanks{Liang~Zhao is with the School of Computer Science, Shenyang Aerospace
 University, Shenyang 110135, China (e-mail: lzhao@sau.edu.cn).}
}

% The paper headers
\markboth{J. of \LaTeX\ C. F.,~Vol.~14, No.~8, Aug.~2021}%
{Shell \MakeLowercase{\textit{et al.}}: A Sample Article Using IEEEtran.cls for IEEE Journals}

% \IEEEpubid{0000--0000/00\$00.00~\copyright~2021 IEEE}
% Remember, if you use this you must call \IEEEpubidadjcol in the second
% column for its text to clear the IEEEpubid mark.

\maketitle

\begin{abstract}
\rev{The rapid electrification of modern transportation necessitates degradation modeling methods that provide remaining useful life (RUL) prediction and state of health (SOH) estimation for vehicular battery management systems.}
Robustness to out-of-distribution (OOD) degradation patterns driven by internal and external factors, \rev{like materials and dynamic vehicular operating conditions, is a critical requirement for these on-board monitoring methods.}
% 
However, achieving such robust methods capable of handling the prevalent distribution shifts presents significant challenges, e.g., spatiotemporal information scarcity, domain shift-induced feature statistics uncertainty, and early-stage bias in multi-source imbalanced degradation distribution.
To address these issues, we propose a robust OOD degradation tracking network (RODTNet) and its long-tail optimized variant (RODTNet-LT) tailored \rev{for electric vehicle powertrain monitoring} via uncertainty-driven feature manipulation with topological constraint.
% 
Firstly, we design a manifold-contraction-based degradation representation extraction strategy using an intrinsic-dimension-reducing and degradation-continuity-preserving signal reconstruction and a Mamba-based encoder.
%
Secondly, we propose an uncertainty-driven feature synthesis method featuring probabilistic degradation characterization and quantifying statistics uncertainty via a multivariate Gaussian distribution to mitigate domain shifts.
%
Thirdly, we develop a long-tail-aware adaptive rebalancing refinement strategy that utilizes neighboring degradation information through Gaussian kernel smoothing to amplify underrepresented sample representations and address early-stage bias.
%
Finally, we establish 11 cross-domain generalization tasks across 4 datasets including 386 cells, with benchmarking against 14 baselines (with 20 variants) validating the superior robustness of both RODTNet and RODTNet-LT, \rev{proving their suitability for reliable vehicular electronic systems.}
\end{abstract}

\begin{IEEEkeywords}
\rev{Automotive Li-ion battery},
\rev{vehicular battery management systems},
degradation modeling,
domain generalization,
feature manipulation,
long-tail rebalancing.
\end{IEEEkeywords}

% \bstctlcite{IEEEexample:BS}

\section{Introduction}

% para 1
\rev{As electric vehicle become the cornerstone of sustainable mobility, Li-ion batteries (LIBs) serve as the fundamental energy storage units within vehicular powertrains\cite{10766661, 9849057, 10897846, WASEEM2025236414}.} However, degradation through cycling compromises driving range, elevates thermal runaway risks, and diminishes residual value for second-life utilization \cite{LI2024155588}.
\rev{These critical challenges underscore the imperative for precise degradation monitoring through a key pair of complementary indicators.} Namely, remaining useful life (RUL) which provides a forward-looking prediction by quantifying the operational cycles left until capacity drops below the end-of-life threshold, and state of health (SOH) which reflects the current condition via the current-to-nominal capacity ratio.

% para 2
\rev{Due to the requirements for non-invasive measurements in on-board applications, RUL and SOH of batteries can only be inferred indirectly through degradation modeling leveraging monitoring data acquired from the vehicular battery management systems (BMS).}
% 
While traditional mechanistic approaches struggle with complex multi-physics coupling effects, data-driven methods, particularly deep learning based ones, have emerged as more powerful tools for establishing nonlinear mappings between operational signals and degradation states\cite{XIAO2023107347, XIE2024630}.
However, their performance heavily relies on the assumption that training and testing data share similar distributions, \rev{which is violated in dynamic real-world driving conditions.} Furthermore, the prohibitive cost of acquiring labeled data across all operational scenarios compounds the generalization challenges.

% para 3
Numerous studies have investigated training-testing distribution shifts in battery systems and proposed models via pretraining-finetuning or domain adaptation \cite{10177255}. \rev{However, these methods remain constrained by requiring target domain data, which are unavailable in deployed vehicular systems due to data transmission bandwidth limits, privacy regulations, and real-time operational constraints\cite{10886998, 10663939, 9849057}.}
% 
Degradation modeling fundamentally differs from discrete-label tasks like anomaly detection and fault diagnosis \cite{10214203, 10577994} due to its continuous target space, where target values exhibit an intrinsic order. This ordinality renders standard approaches inadequate and makes out-of-distribution (OOD) generalization using only source data exceptionally challenging.
%
While some approaches leveraging representation learning or meta learning enable target-free modeling \cite{TAN2024103725, CHEN2024234696}, they still rely on degradation priors. In contrast, data manipulation, a distinct domain generalization paradigm \cite{9782500}, offers advantages in structural simplicity and model-agnostic integration. However, this strategy faces three challenges related to multi-stage propagation from raw inputs to the degradation targets.

\begin{figure*}[t!]
    \vspace{-1em}
	\centering
	\includegraphics[width=7.1in]{Fig1_Challenges.eps}
	\caption{Challenges in OOD-robust degradation modeling of LIBs, including: (a) Spatiotemporal Information Scarcity, (b) Domain Shifts-induced Statistics Uncertainty, and (c) Long-tail Degradation Distribution.}
	\label{fig:challenges}
\end{figure*}

% para 4
The first challenge is the spatiotemporal information scarcity, fundamentally stemming from the non-invasion measurement requirements.
% 
\rev{Temporally, equal-interval sampling results in sparse data within dynamically complex, degradation-sensitive periods \cite{9648002, BOSE2023107543, zhu2022data, GUAN2018204, XIONG2020110048}, while over-representing less informative static operational segments. This leads to an insufficient representation of critical temporal dynamics.}
% 
Spatially, cyclic data relies on a limited set of sensor modalities, providing restricted multi-physical insight and failing to adequately capture crucial internal battery behaviors like thermodynamic stresses \cite{XIONG2020110048} and mechanical deformations \cite{zeng2023nonintrusive}.
% 
As illustrated in Fig.~\ref{fig:challenges}(a), direct modeling of such sparse raw cyclic signals yields a complex high-dimensional manifold structure where degradation representations lack topological continuity. Consequently, subsequent feature manipulation becomes susceptible to noise overfitting and demonstrates limited generalization capacity beyond training distributions.

% para 5
The second challenge arises from the lack of consideration of statistics uncertainty of representations within feature space, induced by domain shifts.
% 
The inherent complexity and sensitivity of battery aging \cite{LI2024155588} cause degradation patterns to inevitably diverge across different operational conditions. This divergence creates domain shifts, introducing statistics uncertainty as source-domain-trained models encounter OOD data with different underlying distributions \cite{9577663, pmlr-v202-dai23d, liuncertainty}.
% 
As illustrated in Fig.~\ref{fig:challenges}(b), these shifts primarily involve covariate shifts caused by variations in charging protocols, electrode materials, and depth of discharge, coupled with characteristic concept shifts arising from divergent degradation trajectories. 
% 
Consequently, this domain shift-induced uncertainty disrupts the consistency of the mapping learned between features and degradation states, limiting a model's ability to generalize from the source domains to unseen target domains and hindering the practical applicability of data-driven prognostics.

% para 6
The third challenge is the multi-domain long-tail degradation distribution, arising from complex internal thermoelectric coupled reactions and appearing as diverse nonlinear trajectories.
%
Specifically, the multi-stage dynamics of battery aging, often marked by a `knee point' transitioning from slow early degradation to accelerated late-stage deterioration, inherently leads to cyclic degradation samples being heavily concentrated in the early-life phase, with later stages sparsely represented.
% 
As illustrated by the fast charging example in Fig.~\ref{fig:challenges}(c), this creates an asymmetric, long-tailed sample distribution. Consequently, models neglecting this characteristic exhibit poor generalization on the less frequent `tail' samples. 
% 
This long-tail distribution biases feature manipulation-based approaches towards the abundant early-life data, creating skewed representations that undermine generalizability. Furthermore, the heterogeneity of these long-tail distributions across multiple source domains exacerbates feature space fragmentation, impeding robust degradation modeling in practice.

% para 7
In this paper, we propose a robust OOD degradation tracking network (RODTNet) and its long-tail-optimized variant (RODTNet-LT) to address these challenges.
% 
First, we propose a manifold contraction-based strategy to extract robust degradation representations (MCRDR), aimed at overcoming spatiotemporal information scarcity. This strategy analyzes feature space manifold topology to identify geometric properties of well-generalized representations, and then utilizes a heuristic signal reconstruction method in conjunction with the Mamba mechanism for feature mapping.
% 
Next, we develop an uncertainty-driven feature synthesis (UDFS) method to cope with domain shifts. By quantifying and sampling from the inherent statistics distribution of target domains via multivariate Gaussian distribution, the UDFS expands the boundaries of the learned degradation manifold and facilitates domain generalization for degradation patterns.
% 
Finally, we propose a long-tail-aware adaptive rebalancing refinement (LARR) method to address the generalization impairment caused by multi-domain long-tail distributions. Using Gaussian kernel smoothing, this method guides the feature synthesis process towards exploring low-density degradation patterns, thereby promoting a more balanced representation within the label space.
% 
The contributions of this work are as follows.

% para 8, 9, 10, 11
\begin{enumerate}
	\item We propose the manifold contraction-based strategy for robust degradation representation extraction (MCRDR). By enforcing pattern continuity and contracting the feature manifold, MCRDR yields a robust degradation representation to overcome the spatiotemporal information scarcity and promotes generalization.
	\item We develop the uncertainty-driven feature synthesis (UDFS) method to tackle domain shifts. UDFS explicitly links feature space statistics uncertainty to domain variability, thereby guiding the extraction of domain-invariant features and the synthesis of diverse OOD patterns for improved generalization.
	\item We present the long-tail-aware adaptive rebalancing refinement (LARR) to address the multi-domain long-tail challenge. Leveraging similarities between feature and label spaces and distribution smoothing, LARR mitigates head-bias of feature synthesis, enhancing generalization for imbalanced degradation trajectories.
	\item Comprehensive experiments across 11 OOD generalization tasks using 4 LIB aging datasets validate the effectiveness of both RODTNet and RODTNet-LT, \rev{and demonstrate their superiority over 14 baseline methods with their 20 variants, advancing reliable vehicular BMS.}
\end{enumerate}

% para 12
This paper is structured as follows.
In Section II, the methods related to robust prognosis are reviewed.
In Section III, the proposed models, RODTNet and RODTNet-LT, for robust degradation modeling are detailed.
In Section IV, the comprehensive evaluations of RODTNet and RODTNet-LT are conducted and the experimental results are presented.
In Section V, this study is concluded.

\section{Related Work}

In this section, we review methods related to robust battery degradation modeling, including methods based on (1) pretraining-finetuning paradigm, (2) domain adaptation paradigm, and (3) domain generalization paradigm.
% (1) pretraining-finetuning requiring full target-domain labels,
% (2) domain adaptation aligning feature distributions with unlabeled target data,
% (3) domain generalization operating without target-domain supervision.

\subsection{Pretraining-finetuning-based Degradation Modeling}

In \cite{LI2021116410}, Li \emph{et al.} propose a lightweight and generalizable model for battery SOH estimation that integrates the pretraining-finetuning paradigm with pruned convolutional neural network.
\rev{In \cite{9372902}, Bhattacharjee \emph{et al.} utilize finetune 1D CNN to facilitate state estimation for varying battery chemistries.}
\rev{In \cite{9343713}, Che \emph{et al.} introduce an online self-calibrating model for RUL prediction based on Gaussian process regression, long short-term memory neural network, and a pretraining-finetuning strategy.}
In \cite{LIU2025115347}, Liu \emph{et al.} fine-tune a bidirectional long short-term memory (BiLSTM)-Mamba2 model based on differential voltage features to achieve SOH estimation across batteries with different operating conditions and electrode materials.
In \cite{D2EE01676A}, Ma \emph{et al.} propose a real-time personalized battery degradation modeling framework named CRNN, based on convolutional-recurrent neural network, achieving joint prediction of SOH and RUL across operating conditions.
In \cite{wang2024physics}, Wang \emph{et al.} propose a physics-informed neural network (PINN) for battery degradation modeling and achieve cross-dataset battery health state prediction using the pretraining-finetuning paradigm.
\rev{In \cite{10886998}, Wang \emph{et al.} propose a decentralized framework to achieve robust SOH estimation across heterogeneous battery datasets.}

\subsection{Domain Adaptation-based Degradation Modeling}

In \cite{9788003}, Ma \emph{et al.} propose a method named ConvMMD that aligns source and target domains by minimizing the maximum mean discrepancy of lower-order moments between them, facilitating a condition-robust estimation of SOH of batteries.
In \cite{HAN2022230823}, Han \emph{et al.} present a high-precision degradation model requiring only a few target domain samples, achieved by incorporating a domain adaptation layer into long short-term memory neural networks, known as LSTM-DA.  
In \cite{9560040}, Ye \emph{et al.} design a method named DDAN that realizes unsupervised degradation feature alignment via adversarial training across source and target domains, enabling a transferable prognostic model.
In \cite{10029904}, Lu \emph{et al.} develop a two-stage prognostic framework named DKTL, which combines degradation mechanisms and achieves knowledge transfer of degradation patterns from synthetic data to real data.
In \cite{WANG2023108897}, Wang \emph{et al.} develope a dual-branch model with a domain invariant-specific feature decoupling strategy and a degradation trend preservation strategy to align degradation patterns across domains, known as DR-Net.

\subsection{Domain Generalization-based Degradation Modeling}

Domain generalization paradigm involves three main categories, namely data manipulation, representation learning, and specific learning strategy \cite{9782500}.
In \cite{CHEN2024234696}, Chen \emph{et al.} develop a dual-branch feature extractor to process domain samples with mixed and separated sources. By leveraging a domain discriminator and a health state estimator, their approach enables domain-invariant degradation feature extraction, culminating in a representation learning-based degradation pattern generalization method termed DGSOH.
In \cite{TAN2024103725}, Tan \emph{et al.} design a meta learning-based domain generalization method named MAGNet. It couples a two-phase meta-optimization strategy (training/testing cycles) with RUL-driven health prediction, enforcing domain-invariant representation learning through task-aligned gradient synchronization. % promoting the learning of cross-domain consistent representations.
In \cite{liu2024itransformer}, Liu \emph{et al.} propose a method designed to obtain robust sequence representations against OOD temporal patterns through the inclusion of a transposed attention module and a feed-forward network module, known as iTransformer.
In \cite{nie2023a}, Nie \emph{et al.} present a time-series modeling framework resistant to pattern drift known as PatchTST. The model employs subsequence segmentation and channel-independent modeling to enhance its adaptability to dynamic temporal patterns.
In \cite{chen2023tsmixer}, Chen \emph{et al.} develop a purely linear model capable of efficiently extracting spatio-temporal dependencies without relying on task-specific prior knowledge, known as TSMixer.
% 
\rev{In \cite{10980357}, Alamin \emph{et al.} propose a digital twin framework to realize effective cloud-vehicle generalization for BMS.}

While pretraining-finetuning and domain adaptation paradigms exhibit varying levels of dependency on target domain data, approaches employing these strategies face deployment constraints in scenarios with dynamic operating conditions and partial observability.
% Representation learning and meta-learning strategies eliminate the need for target domain samples but often rely on precise domain labels or complex pretext tasks, limiting their effectiveness in intricate transportation environments.
While the representation learning-based and meta learning-based generalization methods enable target-free modeling, their dependence on degradation priors, including auxiliary metadata or pretext tasks, limits scalability and generalization performance.
% 
% These limitations inspire the exploration of methods driven by the relatively understudied data manipulation strategy, particularly the feature manipulation strategy.
Though promising, data manipulation-based generalization method faces open issues in robust battery prognosis, including scarce spatiotemporal degradation information, statistics uncertainty in feature spaces due to domain shifts, and the long-tail degradation distributions across multi-source domains.
To address these issues, we propose RODTNet, a novel model designed to achieve OOD degradation pattern generalization. RODTNet generates continuous degradation patterns within a low-complexity latent feature space by employing heuristic feature reconstruction and probabilistic feature statistics generation. Furthermore, we propose RODTNet-LT, an enhanced variant of RODTNet, which incorporates statistical feature correction to specifically address the long-tail distribution challenge.

\section{Methodology}

\begin{figure*}[t!]
    \vspace{-1em}
	\centering
	\includegraphics[width=7.1in]{Fig2_RODTNet_framework.eps}
	\caption{Overview of the proposed RODTNet and its long-tail optimized variant, RODTNet-LT. The core RODTNet comprises a manifold contraction-based robust degradation representation extraction mechanism (with a Mamba-based encoder) and an uncertainty-driven feature synthesis module. RODTNet-LT extends RODTNet by incorporating a long-tail-aware adaptive rebalancing refinement module.}
	\label{fig:framework}
\end{figure*}

\subsection{Overview}

\rev{To advance reliable prognostic capabilities in intelligent vehicular BMS and enhance system robustness against unseen degradation patterns, we propose RODTNet, a robust OOD degradation tracking network, and its long-tail optimized variant, RODTNet-LT.} As illustrated in Fig.~\ref{fig:framework}, the core RODTNet integrates two key components: (1) a manifold contraction-based robust degradation representation extraction (MCRDR) strategy, and (2) an uncertainty-driven feature synthesis (UDFS) method. RODTNet-LT further incorporates (3) a long-tail-aware adaptive rebalancing refinement (LARR) module to specifically address challenges in imbalanced degradation distributions.
% 
First, the MCRDR module processes on-board sensor signals via heuristic reconstruction and employs a Mamba model to map them into a latent representation that forms a manifold with low topological complexity and enhanced continuity, addressing the issue of insufficient spatiotemporal coverage.
%
Then, the UDFS module models the feature distribution of seen domain patterns, quantifies the uncertainty arising from domain shifts, and synthesizes plausible features corresponding to unseen domains by leveraging this uncertainty measure, promoting OOD generalization. 
% 
Finally, the LARR module complements the previous stages by adaptively recalibrating and rebalancing the feature space. This refinement sharpens the model's sensitivity to both frequent and rare degradation patterns, addressing the long-tail distribution challenge.

We formulate the problem of modeling battery degradation under distribution shifts. \rev{Let $\boldsymbol{x}_{t} \in \mathbb{R}^{L \times C}$ denote the multi-sensor monitoring signal sequence (length $L$, $C$ sensors) at cycle $t$, and $y_{t} \in \mathbb{R}$ be the corresponding degradation indicator (RUL or SOH).} We consider $(\boldsymbol{x}_{t}, y_{t})$ as samples from random variables $(X, Y)$ residing in spaces $\mathcal{X} \times \mathcal{Y}$, governed by a joint distribution $P_{XY}$ with marginals $P_{X}$ and $P_{Y}$.
% 
\rev{Variations in automotive operating conditions (e.g., driving profiles, ambient temperatures) and internal battery states create distinct operational domains. We consider $K = M + N$ domains, denoted as $D^{i} = \{ (\boldsymbol{x}_{t}, y_{t}) \} \sim P_{XY}^{i}, i = 1, \dots, K$. Crucially, distinct domains exhibit distribution shifts, meaning $P^{i}_{XY} \neq P^{j}_{XY}$ for $i \neq j$.}
% 
The objective of OOD-robust degradation modeling is to learn a  model $f_{\theta}: \mathcal{S} \to \mathcal{T}$, parametered by $\theta$, using data from $M$ source domains $\mathcal{S} = \{ D^{i} \}_{i=1}^{M}$.
% 
The model should generalize to $N$ unseen target domains $\mathcal{T} = \{ D^{j} \}_{j=M+1}^{M+N}$ without requiring access to them during training. This is achieved by minimizing the expected loss $\ell$ on the target domains, that is $\mathop{\min}_{\theta}\mathbb{E}_{(\boldsymbol{x}, y) \in \mathcal{T}} [\ell(f_{\theta}(\boldsymbol{x}), y)]$ where the expectation is taken over the data distribution of the target domain.

\subsection{Manifold Contraction for Robust Degradation Representation}

To address the challenge of insufficient spatiotemporal coverage, we design MCRDR based on manifold contraction that emphasizes critical temporal segments and achieves spatial dimension expansion, enabling subsequent feature synthesis to further increase feature density.
First, we investigate the relationship between the generalization ability of degradation features and their manifold topology, exploring the geometric properties of well-generalized representations.
Second, we propose a heuristic cyclic signal reconstruction method possessing potentially low topological complexity.
Finally, the reconstructed input are mapped onto the feature space using a Mamba-based model.

\subsubsection{Degenerate Representation from a Topological Perspective}

Inspired by studies in relevant fields such as fault diagnosis, we investigate the geometric properties of generalizable features of degradation models.
Omitting cycle and domain indices for notational clarity, we explicitly consider the latent variable $\boldsymbol{z}$ corresponding to input $\boldsymbol{x}$. Consequently, the domain is defined as $D = \{ (\boldsymbol{x}, \boldsymbol{z}, y) \}$, adapting the source $\mathcal{S}$ and target $\mathcal{T}$ domain structure accordingly.
Let $Z$ be the random variable for $z \in \mathcal{Z}$ with distribution $P_{Z}$, we assume these representations reside on a manifold $\mathcal{M}_{Z}$ embedded within the feature space.
We assess degradation representations using the Information Bottleneck (IB) \cite{alemi2017deep} principle for quality and Intrinsic Dimension (IntDim) \cite{10.5555/3692070.3694510} for geometry.
The IB principle seeks an optimal representation $Z$ by maximizing its relevance to the label $Y$, $\mathcal{I}(Z;Y)$, while minimizing the information it retains about the input $X$, $\mathcal{I}(Z;Y)$. This objective is formulated as
\begin{equation}
    \min_{Z} \mathcal{IB} = \min_{Z} [- \mathcal{I}(Z;Y) + \beta \mathcal{I}(Z;X)], 
    \label{IB_mutual_info}
\end{equation}
where $\mathcal{I}(\cdot;\cdot)$ is mutual information, and $\beta > 0$ is a Lagrangian multiplier. 
The IntDim is the minimum dimension for isometric manifold embedding.
Given the covariate shift inherent in degradation generalization, the IB objective is reformulated, focusing on the relationship between $Z$ and $Y$, formulated as
\begin{equation}
    \min_{Z} \mathcal{IB} = \min_{Z} [(1-\beta) \mathcal{H}(Y \mid Z) + \beta \mathcal{H}(Z \mid Y)], 
    \label{IB_cond_entropy}
\end{equation}
with conditional entropy $\mathcal{H}(\cdot \mid \cdot)$ and weight $\beta \in (0, 1)$.
The $\mathcal{H}(Y \mid Z)$ term encourages label predictability.
The $\mathcal{H}(Z \mid Y)$ term, beyond potential overfitting implications, crucially links the representation to the data manifold's topology.

The generalization error of prognostic models is upper-bounded by $\mathcal{H}(Z \mid Y)$, formulated as
\begin{align}
    & \mathbb{E}_{\{\mathbf{x},\mathbf{z},y\}\thicksim \mathcal{S}}[||f(\mathbf{z})-{y}||_2] -\mathbb{E}_{\{\mathbf{x},\mathbf{z},y\}\thicksim \mathcal{T}}(||f(\mathbf{z})-{y}||_2) \notag \\
    & \leq 2KQ(\mathcal{H}(Z \mid Y)), \label{generalization_error}
\end{align}
% $||\cdot||_2$ is the $\ell_{2}$ norm
where $K$ is down-task regressor's Lipschitz constant, and $Q(\cdot)$ depends on $P(Z \mid Y)$, influencing the bound's tightness.
Furthermore, $\mathcal{H}(Z \mid Y)$ itself is bounded by the average IntDim of the conditional manifolds $\mathcal{M}_i\subset\mathcal{M}$ corresponding to observations $\boldsymbol{z}$ conditioned on $y = y_{i}$, formulated as
\begin{equation}
    \mathcal{H}(Z \mid Y) \leq \mathbb{E}_{\mathbf{y}_{i} \sim \mathcal{Y}}[-\log(\epsilon)IntDim(\mathcal{M}_i)+\log\frac{K^{\prime}}{C(\epsilon)}],
    \label{intrinsic_dimension} 
\end{equation}
where $\epsilon \to 0^{+}$ is a neighborhood radius, $K^{\prime}$ is a constant, and $C(\epsilon)$ is some function of $\epsilon$. 
% $i$ indexes samples
% $ID(\cdot)$ is the intrinsic dimension of the manifold
Conditions, proof methodologies and other details for these bounds \cite{Ghosh_Motani_2023, 10.5555/3692070.3694510} are satisfied by our model, and details are omitted for brevity.

Based on Eqs.~\eqref{generalization_error} and \eqref{intrinsic_dimension}, as well as the IB principle's balance between compression and prediction, desirable degradation representations should possess two key geometric properties. These are: 
a low IntDim, reflecting a relatively simple manifold structure faithfully representable in lower dimensions, and continuous degradation patterns, represented as ordered and separable feature distributions for different aging states.
From a domain generalization perspective, these properties are vital. First, a low IntDim facilitates manifold contraction, simplifying the representation structure within which feature synthesis and distribution extrapolation are performed.
Second, continuity and order across degradation states, reflected in separable feature distributions, promote domain invariance and enhance feature manipulation by reducing inter-state interference and enabling reliable density estimation.

% Feature Engineering as Manifold Contraction
Acquiring topologically optimal degradation representation via IntDim regularization on raw signals is challenging due to costly per-step IntDim estimation and feature discontinuities from sparse coverage, impeding effective optimization. Consequently, inspired by PINN methods in battery research \cite{wang2024physics, 10251604}, we utilize a heuristic cycle signal reconstruction.
We implement this reconstruction by extracting statistical features from each sensor sequence of the original signal $\boldsymbol{x}_{t} \in \mathbb{R} ^ {L \times C}$. 
Representative examples include temporal differences or channel-wise moments such as mean and standard deviation. Concatenating these extracted statistics yields the reconstructed signal $\boldsymbol{\widetilde{x}}_{t} \in \mathbb{R} ^ {D \times 1}$, where $D$, the total number of extracted statistics, typically satisfies $C < D \ll L$. The precise input construction methodology is detailed in Section~\ref{subsec:experimental_setup}.
Empirically validated in Section~\ref{subsec:comparative_experiments}, this reconstruction mitigates temporal interference and expands features dimensionality to address insufficient coverage challenge, resulting in lower IntDim and more continuous degradation representation.

\subsubsection{Mamba with Manifold Contraction}

\rev{Selecting a backbone for mapping reconstructed signals to degradation features involves considering the task's sequential nature, linear time-invariant (LTI) system approximability for battery physics, and vehicular edge deployment constraints. Mamba's sequence modeling proficiency, LTI-compatible state-space mechanism, linear scaling, and parallel inference capabilities \cite{gu2024mamba} align well with these requirements, leading us to adopt it for this mapping task.}

We form Mamba's input $\widetilde{X}_{Recon} \in \mathbb{R}^{D \times E}$ by vertically stack $E$ consecutive reconstructed signals.
This input is first embedded channel-independently into $\widetilde{X}_{Emb} \in \mathbb{R}^{H \times E}$, where $H$ is the hidden state dimension, then mapped via a three-branch residual structure into the feature space as
\begin{align}
& Z^{B}_{MRB} = SiLU(Linear(\widetilde{X}_{Emb})), \\
& Z^{M}_{MRB} = SSM(SiLU(Conv(Linear(\widetilde{X}_{Emb})))), \\
& Z_{MRB} = Linear(Z^{M}_{MRB} \odot Z^{B}_{MRB}) + \widetilde{X}_{Emb},
\end{align}
where $Z^{M}_{MRB}$, $Z^{B}_{MRB}$, and $Z_{MRB}$ are the main, bypass, and final Mamba residual block output. Operations include concatenation $\odot$, linear mapping $Linear(\cdot)$, convolution $Conv(\cdot)$, SiLU activation $SiLU(\cdot)$, and the S6 transformation $SSM(\cdot)$. The S6 model, an input-dependent discrete recurrent neural network, is formulated as
\begin{align}
& h_{t} = \overline{A}h_{t-1} + \overline{B}\widetilde{x}_{t}, \\
& \overline{A} = DR_{A}(\Delta, A) = exp(\Delta A), \\
& \overline{B} = DR_{B}(\Delta, A, B) = (\Delta A)^{-1}(exp(\Delta A)-I)\cdot\Delta B,
\end{align}
where $A$ and $B$ are trainable matrices. Adopting ZOH discretization $DR$ \cite{gu2024mamba}, $A$ is HiPPO-initialized \cite{20212610553629}, while $B$ and the discretization factor $\Delta$ are functions of input $\widetilde{x}_{t}$.

\subsection{Uncertainty-driven Feature Synthesis against Domain Shifts}\label{subsec:UDFS}

To mitigate the issue of domain shift-induced statistics uncertainty and enhance model generalization capability, we propose UDFS method, which extrapolates distribution boundaries based on the contracted degradation representation.
First, we model feature statistics probabilistically via multivariate Gaussian distributions.
Second, we quantify statistics uncertainty by non-parametrically estimating the variance of these statistics.
Lastly, based on this uncertainty quantification, we synthesize features representing potential OOD degradation patterns through sampling, and the UDFS mechanism is subsequently integrated into the MMC model to enhance robustness against OOD instances.

\subsubsection{Multivariate Gaussian Distribution for Domain Shift Modeling}

Domain-specific degradation patterns are reflected in statistical distributions of features. We posit that domain shifts induce quantifiable inconsistencies in these statistics across different operational environments. Explicitly modeling this statistical inconsistency, particularly through the uncertainty inherent in feature statistics, enables the synthesis of features representative of unseen target domains from source data. This approach allows training degradation models with enhanced robustness.
Our method quantifies inter-domain variability by identifying the underlying distribution that governs feature statistics and by leveraging its sampling variance.
% 
Quantifying statistics uncertainty for robust degradation modeling essentially involves characterizing the low-order moments of feature statistic distributions.
This requires specific properties from the underlying distribution.
Ideally, it should satisfy Carleman's condition \cite{akhiezer2020classical}, ensuring unique determination by its moments, while its low-order moments should capture suffici distributional information. The multivariate Gaussian distribution, which satisfies Carleman's condition and is completely defined by its mean and covariance, fulfills these requirements. Given this theoretical alignment and its established use in battery degradation analyses, we employ a multivariate Gaussian framework to model feature statistics.

Specifically, for features $ Z_{MRB} \in \mathbb{R}^{B \times H \times E} $ derived from a minibatch of $ B $ samples, we employ batch-wise mean $ \mu_{Z} $ and standard deviation $ \sigma_{Z} $ as their statistical representations.
As described in the previous paragraph, we have $ \mu_{Z} \sim \mathcal{N}(\mu, \Sigma_{\mu}^{2}) $ and $ \sigma_{Z} \sim \mathcal{N}(\sigma, \Sigma_{\sigma}^{2}) $, where $ \mu $ and $ \sigma $ are the means of the feature statistics $ \mu_{Z} $ and $ \sigma_{Z} $, respectively.
We compute these statistics directly using the mean and standard deviation of features, formulated as
\begin{align}
    & \mu=\frac{1}{HE}\sum_{H}\sum_{E}Z_{MRB}, \\
    & \sigma = \sqrt{\frac{1}{HE}\sum_{H}\sum_{E}(Z_{MRB} - \mu)^{2}}.
\end{align}
And $ \Sigma_{\mu}^{2} $ and $ \Sigma_{\sigma}^{2} $ are the variances associated with feature statistics $\mu_{Z}$ and $\sigma_{Z}$, which also serve to model the statistics uncertainty induced by domain shifts.
These variances are estimated non-parametrically based on the observed intra-batch variability of the feature statistics, formulated as
\begin{align}
    & \Sigma_{\mu}^{2}=\frac{1}{B}\sum_{b\in B}(\mu-\mathbb{E}_{b}(\mu))^{2}, \\
    & \Sigma_{\sigma}^{2} = \frac{1}{B}\sum_{b\in B}(\sigma-\mathbb{E}_{b}(\sigma))^{2},
\end{align}
where $ \mathbb{E}_{b}(\mu) $ and $ \mathbb{E}_{b}(\sigma) $ are the expectations of mean and standard deviation of the $ b $-th sample within a minibatch.

\subsubsection{Uncertainty-driven Feature Synthesis}

Differentiable generation of feature statistics employs the reparameterization method under the assumption of Gaussian distributed original statistics with parameters $\mu$, $\Sigma_{\mu}$ and $\sigma$, $\Sigma_{\sigma}$, formulated as
\begin{align}
    & \mu_{Z}^{\prime}=\epsilon_{\mu}\Sigma_{\mu}+\mu,  \\
    & \sigma_{Z}^{\prime}=\epsilon_{\sigma}\Sigma_{\sigma}+\sigma, 
\end{align}
where $\mu_{Z}^{\prime}$ and $\sigma_{Z}^{\prime}$ are the generated feature mean and standard deviation, and $\epsilon_{\mu}$ and $\epsilon_{\sigma}$ are samples drawn from a standard normal distribution $\mathcal{N}(0, 1)$.
% 
The proposed UDFS utilizes the mean statistic $\mu$ as an anchor and standard deviation $\Sigma_{\mu}$ and $\Sigma_{\sigma}$ to probabilistically parameterize a offset, conferring adaptability to stochastic domain shifts through its variable direction and magnitude. Generating new statistics via reparameterization introduces controlled randomness, which augments the feature manifold surrounding the anchor statistics. This expansion improves coverage of potential unseen degradation patterns, enhancing model robustness and enabling generalization to OOD scenarios.
% 
Functionally, UDFS is integrated into the Mamba with manifold contraction model, operating between the encoder and the task-specific regression head for feature synthesis, formulated as
\begin{equation}
    G(Z_{MRB}) = \sigma_{Z}^{\prime}\left(\frac{Z_{MRB} - \mu}{\sigma}\right) + \mu_{Z}^{\prime}.
    \label{UDFS}
\end{equation}
This UDFS integration adds negligible constant-level computational overhead during training.
Furthermore, a dynamic activation mechanism based on Bernoulli sampling regulates its use, ensuring sufficient learning from real samples. At each training iteration, a Bernoulli value $p^{\prime}$ is drawn, and feature synthesis via UDFS proceeds only when $p^{\prime} < p_{0}$, where $p_{0}$ is a predefined activation threshold.

\subsection{Long-tail-aware Adaptive Rebalancing Refinement}\label{susbec:LARR}

To address the challenge of multi-source domain long-tail degradation distributions, we propose LARR, which employs feature distribution calibration specifically for enhancing battery degradation trajectory tracking and precise SOH estimation.
First, we employ feature distribution calibration to strengthen the correspondence between degradation feature similarity and degradation label proximity.
Second, we integrate this feature calibration strategy into UDFS to form LARR to eliminate domain-specific information tied to long-tail distributions, thereby purifying feature domain invariance.
Finally, we examine LARR's training dynamics and subsequently incorporate a head-bias mitigation strategy designed to improve degradation pattern generalization.

\subsubsection{Symmetric-kernel-based Feature Statistic Calibration}

We discretize the bounded SOH label space via uniform binning into $\mathcal{B}$ intervals, indexed $\textit{b} = 1, 2, \dots, \mathcal{B}$. Battery degradation trajectories exhibit a long-tail distribution, characterized by substantial, often exponential, variation in the number of samples associated with each bin $\textit{b}$.
% 
Continuous degradation labels imply feature similarity relates inversely to label distance, unlike discrete labels. 
% 
The inherent long-tail distribution of degradation data violates this, as feature scarcity in tail bins leads to poor representations spuriously similar to head bin features.
% 
To counteract this, we calibrate features within each bin $\textit{b}$ by propagating information from adjacent bins via symmetric kernel functions. These corrected features, combined with the original observed features from the encoder-UDFS pipeline, are subsequently transformed to achieve a more balanced degradation feature representation.

We compute bin-wise feature statistics similarly to the batch-wise counterparts described in Section \ref{subsec:UDFS}. For features $Z_{MRB}$ within bin $\textit{b}$, the mean $\mu_{\textit{b}}$ and standard deviation $\sigma_{\textit{b}}$ are formulated as
\begin{align}
    & \mu_{\textit{b}} = \frac{1}{HE}\sum_{H}\sum_{E}Z_{MRB}, \\
    & \sigma_{\textit{b}} = \sqrt{\frac{1}{HE}\sum_{H}\sum_{E}(Z_{MRB} - \mu)^{2}}.
\end{align}
% 
To calibrate the feature distribution, we employ kernel smoothing across bins based on their corresponding SOH values. Without loss of generality, we select the Gaussian kernel for this purpose. This yields calibration statistics $\widetilde{\mu}_{\textit{b}}$ and $\widetilde{\sigma}_{\textit{b}}$, formulated as
\begin{align}
    \widetilde{\mu}_{\textit{b}} & = \sum_{1 \leq \textit{b}, \textit{b}^{\prime} \leq \mathcal{B}}Gau(y_{\textit{b}}, y_{\textit{b}^{\prime}})\mu_{\textit{b}^{\prime}} \notag \\
    & = \sum_{1 \leq \textit{b}, \textit{b}^{\prime} \leq \mathcal{B}}exp(-\gamma\|y_{\textit{b}} - y_{\textit{b}^{\prime}}\|_{2})\mu_{\textit{b}^{\prime}},\label{mu_calibration} \\
    \widetilde{\sigma}_{\textit{b}} & = \sum_{1 \leq \textit{b}, \textit{b}^{\prime} \leq \mathcal{B}}Gau(y_{\textit{b}}, y_{\textit{b}^{\prime}})\sigma_{\textit{b}^{\prime}} \notag \\
    & = \sum_{1 \leq \textit{b}, \textit{b}^{\prime} \leq \mathcal{B}}exp(-\gamma\|y_{\textit{b}} - y_{\textit{b}^{\prime}}\|_{2})\sigma_{\textit{b}^{\prime}},\label{sigma_calibration}
\end{align}
where $Gau(\cdot, \cdot)$ is the Gaussian kernel function with coefficient $\gamma$, and $y_{\textit{b}}$ and $y_{\textit{b}^{\prime}}$ are SOH values corresponding to the $\textit{b}$-th and $\textit{b}^{\prime}$-th bins, respectively.

\subsubsection{Imbalance-bias-free Degradation Dynamics Generalization}

To address the long-tail distribution challenge, which persists in a consistent form both preceding and following the feature synthesis via $G(Z_{MRB})$, we implement the proposed LARR algorithm for head-bias-free correction of features.
LRAA performs feature distribution calibration utilizing two sets of bin-specific statistics, namely observed statistics $\{{\mu}_{\textit{b}}, {\sigma}_{\textit{b}}\}$ and calibration statistics $\{\widetilde{\mu}_{\textit{b}}, \widetilde{\sigma}_{\textit{b}}\}$, formulated as
\begin{equation}
    \widetilde{G}(Z_{MRB})=\widetilde{\sigma}_{\textit{b}}\left(\frac{G(Z_{MRB}) - \mu_{\textit{b}}}{\sigma_{\textit{b}}}\right)+\widetilde{\mu}_{\textit{b}}.
\end{equation}
Structurally identical to Eq.~\eqref{UDFS}, the LARR approach, similar to UDFS, is applied exclusively during the training stage. % At iteration , the adjustment is performed via:
To facilitate clear exposition of the training dynamics, we denote the iteration-specific observed statistics and calibration statistics for bin $\textit{b}$ as $\{\mu_{\textit{b}}^{(i)}, \sigma_{\textit{b}}^{(i)}\}$ and $\{\widetilde{\mu}_{\textit{b}}^{(i)}, {\widetilde{\sigma}_{\textit{b}}^{(i)}}\}$, respectively, where $i$ indexes the training iteration. Within each training iteration, the degradation representation is adjusted, and subsequently, the corresponding calibration statistics are updated at the iteration's conclusion.
To ensure training stability, the observed statistics are updated via a momentum-based mechanism that integrates their values from the previous iteration and the static statistics, formulated as
\begin{align}
    & \mu_{\textit{b}}^{(i)} \leftarrow \alpha\mu_{\textit{b}}^{(i-1)}+(1-\alpha)\mu_{\textit{b}}, \\
    & {\sigma_{\textit{b}}}^{(i)} \leftarrow \alpha{\sigma_{\textit{b}}}^{(i-1)}+(1-\alpha)\sigma_{\textit{b}},
\end{align}
where $\alpha$ is the update momentum. The calibration statistics are updated according to Eqs.~\eqref{mu_calibration} and \eqref{sigma_calibration}.

\begin{figure*}[t!]
    \vspace{-1em}
    \centering
    \subfloat{
        \includegraphics[width=1.1in]{fig3_deg_traj_matr_1.eps}
        \includegraphics[width=1.1in]{fig3_deg_traj_matr_2.eps}
        \includegraphics[width=1.1in]{fig3_deg_traj_tongji_1.eps}
        \includegraphics[width=1.1in]{fig3_deg_traj_tongji_2.eps}
    \includegraphics[width=1.1in]{fig3_deg_traj_hust.eps}
        \includegraphics[width=1.1in]{fig3_deg_traj_xjtu.eps}
    }

  \subfloat{
        \includegraphics[width=1.1in]{fig3_tri_lt.eps}
        \includegraphics[width=1.1in]{fig3_mit_lt.eps}
        \includegraphics[width=1.1in]{fig3_tongji_nca_lt.eps}
        \includegraphics[width=1.1in]{fig3_tongji_ncm_lt.eps}
    \includegraphics[width=1.1in]{fig3_hust_lt.eps}
        \includegraphics[width=1.1in]{fig3_xjtu_lt.eps}
    }
    \caption{Illustration of degradation trajectories, domain splits, and the long-tail data distribution of batteries in the MATR, Tongji, HUST, and XJTU dataset.}
    \label{fig:degradation_trajectory}
\end{figure*}

\section{Experiments}

% 数据集描述
\begin{table*}[h!]
    \centering
    \caption{Dataset Overview.}
    \scalebox{0.8}
    {
    \small
    \begin{threeparttable}
      \begin{tabular}{cccccc}
      \toprule
      Dataset & Subset & Cathode materials\tnote{a} & Norminal capacity & Numer of cells & Number of samples \\
      \midrule
      MATR \cite{severson2019data, attia2020closed}  & batch\{1, 2, 3, 9\tnote{b} \} & \ce{LiFePO_{4}} & 1100 mAh  & 144   & 115416 \\
      \midrule
      \makecell{Tongji \cite{zhu2022data}} & \makecell{NCA, NCM\tnote{c}} & \makecell{\ce{Li_{0.86}Ni_{0.86}Co_{0.11}Al_{0.03}O_2(NCA)} \\ \ce{Li_{0.84}Ni_{0.83}Co_{0.11}Mn_{0.07}O_2}(NCM)} & 3500 mAh  & \makecell{116} & \makecell{41963} \\
      \midrule
      HUST \cite{D2EE01676A}  & $\backslash$ & \ce{LiFePO_{4}} & 1100 mAh  & 77    & 114366 \\
      \midrule
      XJTU \cite{wang2024physics} & Batch-\{1, 2, 4\} & \ce{LiNi_{0.5}Co_{0.2}Mn_{0.3}O_2} & 2000 mAh  & 31    & 11182 \\
      \bottomrule
      \end{tabular}%
      \begin{tablenotes}
        \item[a] The anode material in all batteries is graphite.
        \item[b] Only the cells in batch 9 with the protocols `3.6-6-5.6-4.755', `4.8-5.2-5.2-4.16', `6-5.6-4.4-3.834', and `8-4.4-4.4-3.94' are used.
        \item[c] Cell {CY35-05\_1-\#3} from NCA subset and cells {CY35-05\_\#\{1, 2, 3, 4\}} from NCM subset are excluded.
    \end{tablenotes}
  \end{threeparttable}
    }
    \label{tab:dataset}%
  \end{table*}%
  
  % 域泛化任务划分
  \begin{table*}[htbp]
    \vspace{-1em}
    \footnotesize
      \centering
      \caption{The designed OOD degradation pattern generalization tasks.}
      % \scalebox{0.9}
      % {
      %   \small
      \begin{threeparttable}
        % \resizebox{\linewidth}{!}{
        \begin{tabularx}{\textwidth}{ccXX}
        \toprule
        Task & Dataset & ID of cell in SOURCE DOMAIN (\textbf{number of source domain}) & ID of cell in TARGET DOMAIN (\textbf{target domain}) \\
        \midrule
        1     & MATR  & 8C(35\%)-3.6C\_\{47,48\}, 8C(25\%)-3.6C\_\{45,46\}, \newline
        8C(15\%)-3.6C\_\{43,44\}, 7C(40\%)-3C\_\{37,38\}, \newline
        7C(40\%)-3.6C\_\{41,42\}, 7C(30\%)-3.6C\_\{39,40\}, \newline
        6C(50\%)-3.6C\_\{35,36\}, 6C(40\%)-3C\_\{25,26\}, \newline
        6C(40\%)-3.6C\_\{33,34\}, 6C(4\%)-4.75C\_46, \newline
        etc\tnote{a}. (\textbf{52})
        % 6C(30\%)-3.6C\_32, 5C(67\%)-4C\_\{10,18,2,34,42,47\}, \newline
        % 5.9C(60\%)-3.1C\_\{32,40\}, 5.9C(15\%)-4.6C\_\{24,39\}, \newline
        % 5.6C(5\%)-4.75C\_41, 5.6C(38\%)-4.25C\_38, \newline
        % 5.6C(36\%)-4.3C\_\{13,15,21,23,29,37,45,5\}, \newline 
        % 5.4C(70\%)-3C\_\{17,18\}, 5.4C(60\%)-3C\_\{15,16\}, \newline
        % 5.4C(50\%)-3C\_14, 5.4C(40\%)-3.6C\_20, \newline
        % 5.3C(54\%)-4C\_\{11,19,27,35,3,43,48,8\}, \newline
        % 5.2C(58\%)-4C\_34, 5.2C(50\%)-4.25C\_33, \newline
        % 5.2C(37\%)-4.5C\_32, 5.2C(10\%)-4.75C\_31, \newline
        % 4C(80\%)-4C\_\{5,6\}, 4C(40\%)-6C\_29, \newline
        % 4C(4\%)-4.85C\_30, 4.9C(69\%)-4.25C\_26, \newline
        % 4.9C(61\%)-4.5C\_25, 4.65C(44\%)-5C\_22, \newline
        % 4.65C(19\%)-4.85C\_21, 4.4C(80\%)-4.4C\_7, \newline
        % 4.4C(47\%)-5.5C\_18, 4.4C(24\%)-5C\_17, \newline
        % 3.7C(31\%)-5.9C\_\{16,30,38\}, 3.6C(9\%)-5C\_16, \newline
        % 3.6C(80\%)-3.6C\_\{1,2,3\}, 3.6C(30\%)-6C\_15, \newline
        % 5.4C(60\%)-3.6C\_\{23,24\}, 4.4C(8\%)-4.85C\_20, \newline
        % 5.2C(66\%)-3.5C\_35, 3.6C(2\%)-4.85C\_14, \newline
        % 5.6C(19\%)-4.6C\_\{14,20,22,28,36,44,4\}, \newline
        % 4.8C(80\%)-4.8C\_\{17,1,25,9,10,4,7,8,9\}, \newline
        % 4C(13\%)-5C\_27, 5.4C(80\%)-5.4C\_\{11,12\}, \newline
        % 4.65C(69\%)-6C\_23, 6C(60\%)-3C\_\{29,30,48\}, \newline
        % 6C(50\%)-3C\_\{27,28\}, 5.6C(25\%)-4.5C\_37 (\textbf{52})
        & 6C(52\%)-3.5C\_47, 6C(40\%)-4C\_45,
        6C(31\%)-4.25C\_44, \newline
        6C(20\%)-4.5C\_43, 5.6C(65\%)-3C\_42, 5.6C(58\%)-3.5C\_40, \newline
        5.6C(47\%)-4C\_39, 5.2C(71\%)-3C\_36, 4C(31\%)-5C\_28, \newline
        4.9C(27\%)-4.75C\_24, 4.4C(55\%)-6C\_19, 3.6C(22\%)-5.5C\_13, \newline
        2C(7\%)-5.5C\_12, 2C(2\%)-5C\_11, 2C(10\%)-6C\_10, \newline
        1C(4\%)-6C\_9 (\textbf{16}) \\
        \midrule
        2     & MATR  & M-batch9-1\_\{1,10,2,42,20\}, M-batch9-2\_\{21,22,31,36,9\}, \newline M-batch9-3\_\{13,16,23,47,24\} (\textbf{3}) & M-batch9-4\_\{11,12,27,29,38\} (\textbf{1}) \\
        \midrule
        3     & MATR  & M-batch9-4\_\{11,12,27,29,38\}, M-batch9-2\_\{21,22,31,36,9\}, \newline M-batch9-3\_\{13,16,23,47,24\} (\textbf{3}) & M-batch9-1\_\{1,2,10,20,42\} (\textbf{1}) \\
        \midrule
        4     & MATR  & M-batch9-4\_\{11,12,27,29,38\}, M-batch9-1\_\{1,2,10,20,42\}, \newline M-batch9-3\_\{13,16,23,47,24\} (\textbf{3}) & M-batch9-2\_\{21,22,31,36,9\} (\textbf{1}) \\
        \midrule
        5     & MATR  & M-batch9-4\_\{11,12,27,29,38\}, M-batch9-1\_\{1,2,10,20,42\}, \newline M-batch9-2\_\{21,22,31,36,9\} (\textbf{3}) & M-batch9-3\_\{13,16,23,47,24\} (\textbf{1}) \\
        \midrule
        6     & Tongji & CY25-1\_1-\#\{6,1,9,3,7,2,8\}, CY25-05\_1-\#\{14,11,6,\newline 1,2,19,13,7,17,4,16,10,15\},  CY45-05\_1-\#\{28,17,5,2,\newline 21,7,9,19,22,14,16,8,23,24,12,10,20,26,13,27\}, \newline CY25-025\_1-\#\{3,5,4,6\}, CY35-05\_1-\#2 (\textbf{5}) & CY25-1\_1-\#\{4,5\}, CY25-05\_1-\#\{3,5,12,18\}, \newline CY25-025\_1-\#7, CY35-05\_1-\#1, \newline CY45-05\_1-\#\{1,11,15,18,25\} (\textbf{5}) \\
        \midrule
        7     & Tongji & CY25-05\_1-\#\{6,5,13,15,18,7,3,9,1,16,17,4,14,2,12,10,19,11,8\}, \newline CY25-025\_1-\#\{6,7,4,1,2,3,5\}, \newline CY45-05\_1-\#\{15,4,14,28,13,1,3, 26,19,9,18,16,11,12,27,5,10, \newline 22,8,7,21,6,17,24,20,2,25,23\} (\textbf{3}) & CY25-05\_1-\#\{6,23,5,13,15,18,7,3,9,1,\newline 16,17,4,14,21,2,20,12,10,19,11,8,22\}, \newline CY45-05\_1-\#\{23,15,4,14,28,13,1,3,2,25,26,19, \newline 9,18,16,11,12,17,24,20,27,5,10,22,8,7,21,6\} (\textbf{2}) \\
        \midrule
        8     & HUST  & 9-1,2-2,4-7,9-7,1-8,4-6,2-7,8-4,7-2,10-3,2-4,7-4,3-4,5-4,8-7,7-7, \newline 4-4,1-3,7-1,5-2,6-4,9-8,9-5,6-3,10-8,1-6,3-5,2-6,3-8,3-6,4-8,7-8, \newline 5-1,2-8,8-2,1-5,7-3,10-2,5-5,9-2,5-6,1-7,8-3,4-1,4-2,1-4,6-5,4-3, \newline 5-7,3-3,2-3,9-3,10-5,3-2,3-7 (\textbf{55}) & 9-6,4-5,1-2,10-7,1-1,6-1,6-6,9-4,10-4,8-5,5-3, \newline 10-6,2-5,6-2,3-1,8-8,8-1,8-6,7-6,6-8,7-5,10-1 (\textbf{22}) \\
        \midrule
        9     & XJTU  & X-batch2-\{1, $\cdots$\tnote{b}, 15\}, X-batch4-\{1, $\cdots$, 8\} (\textbf{2}) & X-batch1-\{1, $\cdots$, 8\} (\textbf{1}) \\
        \midrule
        10    & XJTU  & X-batch1-\{1, $\cdots$, 8\}, X-batch4-\{1, $\cdots$, 8\} (\textbf{2}) & X-batch2-\{1, $\cdots$, 15\} (\textbf{1}) \\
        \midrule
        11    & XJTU  & X-batch1-\{1, $\cdots$, 8\}, X-batch2-\{1, $\cdots$, 15\} (\textbf{2}) & X-batch4-\{1, $\cdots$, 8\} (\textbf{1}) \\
        \bottomrule
        \end{tabularx}%
        \begin{tablenotes}
          \item[a] Omitted for brevity. Refer to \cite{TAN2024103725} for full details.
          \item[b] This ellipsis mark indicates consecutive cell IDs.
      \end{tablenotes}
      \end{threeparttable}
      \label{tab:tasks}%
      \end{table*}%

\rev{In this section, we present a comprehensive evaluation of the performance of our proposed models, RODTNet and RODTNet-LT, validating their suitability for intelligent vehicular BMS.
We compare them with 14 baseline methods and their 20 variants across 11 generalization tasks derived from 4 datasets encompassing diverse automotive operating conditions.}
We first compare the performance of RODTNet, RODTNet-LT, and baseline methods across diverse datasets and generalization tasks.
Next, we conduct ablation studies to demonstrate the effectiveness of each key component within RODTNet, and to evaluate the specific contribution of the LARR module in RODTNet-LT.
Finally, we investigate the sensitivity of RODTNet's performance to variations in hyperparameter settings.

\subsection{Experimental Setup}\label{subsec:experimental_setup}
  
\subsubsection{Datasets and Generalization Tasks}

We evaluate our framework using 4 large-scale datasets, MATR \cite{severson2019data, attia2020closed}, Tongji \cite{zhu2022data}, HUST \cite{D2EE01676A}, and XJTU \cite{wang2024physics} through 11 OOD generalization tasks. Table~\ref{tab:dataset} summarizes datasets characteristics, including material compositions, electrical specifications, and statistical details, while Fig.~\ref{fig:degradation_trajectory} illustrates the degradation trajectories and the corresponding distribution densities of batteries across different domains from different datasets. Domain division criteria encompass charging/discharging protocols, electrode materials, and depth of discharge. \rev{Specifically, these protocols feature dynamic profiles mimicking real-world electric vehicle driving behaviors.} Tasks 1-5 \cite{TAN2024103725,10522981} systematically examine fast-charging protocol robustness. Task 6 adapts \cite{TAN2024103725}'s configuration to evaluate rare degradation pattern adaptation (0.9\% shared cycles between domains), while Task 7 extends this framework to assess electrode material generalization under OOD conditions. Task 8 \cite{D2EE01676A} evaluates model's robustness to multi-stage discharge protocol variations. Building upon MATR dataset baselines, Tasks 9-11 investigate coupled domain shifts arising from charging protocol-discharge protocol-depth of discharge interactions. Full domain splits appear in Table~\ref{tab:tasks}. Battery degradation is quantified through RUL and SOH, characterizing instantaneous performance and full lifecycle evolution respectively.

\subsubsection{Input Signal Reconstruction}

Following established methodologies in \cite{10251604,wang2024physics}, we implement our heuristic input reconstruction strategy through signal transformations. Dataset-specific reconstruction protocols are implemented to rigorously verify the general applicability of our strategy. For MATR dataset (5 tasks involved), 8-dimensional vectors encode electrochemical descriptors including quadratic coefficients and bias terms derived from differential capacity curve fittings within the voltage range of 2.7–3.3 V, as well as the maximum value, minimum value, variance, cycle-average temperature, internal resistance, and charging time. Other datasets (6 tasks involved) employ 16-dimensional representations to capture multi-phase cycle characteristics. The 16-dimensional feature set is divided into two categories. The first category encompasses statistical metrics, including mean, standard deviation, kurtosis, skewness, curve slope, and curve entropy, extracted from on-board voltage and current sensors within specified operational ranges (current: 0.1–0.5 A; voltage: 4.0–4.2 V for the Tongji and XJTU datasets, and 3.4–3.6 V for the HUST dataset). The second category consists of duration and cumulative charge measurements derived from the constant-current and constant-voltage charging phases.

\subsubsection{Benchmarks}

To evaluate the performance of RODTNet and RODTNet-LT, we conducte a comprehensive comparison across 5 categories of methods, summarized as follows:
% 
(1) state-of-the-art battery prognostic methods, including DeNN \cite{10102497} (with 3 aggregation strategies), PINN \cite{wang2024physics}, and MambaLithium \cite{shi2024mambalithiumselectivestatespace};
(2) pretraining-finetuning-based methods, including CRNN \cite{D2EE01676A} and PINN$^\ast$ \cite{wang2024physics};
(3) domain adaptation-based methods, including ConvMMD \cite{9788003}, LSTM-DA \cite{HAN2022230823}, DTKL \cite{10029904}, DDAN \cite{9560040}, and DR-Net \cite{WANG2023108897};
(4) domain generalization-based methods, including MAGNet \cite{TAN2024103725}, iTransformer \cite{liu2024itransformer}, PatchTST \cite{nie2023a}, and TSMixer \cite{chen2023tsmixer};
(5) RODTNet and RODTNet-LT trained directly on the target domain, referred to as RODTNet$^\ast$ and RODTNet-LT$^\ast$, respectively.
% 
Given the limited application of general-purpose time series forecasting models, including iTransformer, PatchTST, and TSMixer, in the context of battery degradation prediction, we also examine their fine-tuned variants, denoted as iTransformer$^\ast$, PatchTST$^\ast$, and TSMixer$^\ast$.

\subsubsection{Implementation Details}

We conduct all the experiments on a server equipped with an Intel Core i9-13900KF CPU and an NVIDIA 4090 GPU (PyTorch 2.3.1 with CUDA 11.8).
All cycling data undergo min-max normalization to the range between 0 and 1. SOH labels follow dataset-specific scaling protocols, where the MATR dataset utilizes definition-based [0.8-1] scaling while other datasets adopted min-max [0, 1] normalization. RUL labels are normalized during training and denormalized during evaluation.
% 
Details of domain splits are summarized in Table~\ref{tab:tasks}.
% 
Domain generalization-based methods follow a rigorous validation protocol \cite{9847099}, where validation sets are constructed exclusively from source-domain batteries with training-to-validation ratios of 3:1 (MATR batches 1–3, Tongji, XJTU), 4:1 (MATR batch 9), and 6:1 (HUST). All target-domain batteries are reserved for testing. 
% 
To facilitate domain alignment in domain adaptation-based methods, the target domain is further divided into a test subset, which includes batteries designated for degradation modeling, and a training subset comprising the remaining batteries. Both subsets ensure that the labels of degradation samples in the target domain remain completely inaccessible in model training phases.
% 
Models following the pretraining-finetuning paradigm mirror domain generalization protocols during pretraining, utilize identical test sets as domain adaptation methods, and further partition remaining target-domain batteries into 9:1 training-validation splits for finetuning.
% 
Details of sample construction, model architecture, and hyperparameters are provided in Table~\ref{tab:hyperparameters_setting}. The Mean Squared Error loss function in conjunction with the Adam optimizer are used for model parameters optimization. The training process is conducted for a maximum of 100 epochs, with early stopping implemented based on the validation set loss.

% Table generated by Excel2LaTeX from sheet 'Sheet1'
\begin{table}[htbp]
\centering
\caption{Hyperparameter Settings.}
\scalebox{0.8}
{
    \small
\begin{threeparttable}
    \begin{tabular}{rll}
    \toprule
    \multicolumn{1}{l}{Symbol} & Description & Value \\
    \midrule
    \multicolumn{1}{l}{$E$} & Number of consecutive charge-discharge cycles & $1$ \\
    \multicolumn{1}{l}{$H (d_{model})$} & Hidden state dimension & $16$ \\
    \multicolumn{1}{l}{$n_{enc}$} & number of Mmaba residual block & $1$ \\
    \multicolumn{1}{l}{$n_{exp}$} & Expansion factor of Mamba & $4$ \\
    \multicolumn{1}{l}{$r_{\Delta}$} & Rank of discretization factor of Mamba & $1$ \\
    \multicolumn{1}{l}{$d_{convs}$} & Kernel size of CNN in Mamba & $4$ \\
    \multicolumn{1}{l}{$N$} & Dimension of parameter matrices in Mamba & $1024$ \\
    \multicolumn{1}{l}{$p_{0}$} & Activation threshold of UDFA & $0.5$ \\
    \multicolumn{1}{l}{$\mathcal{B}$} & Number of bins of LARR & $100$ \\
    \multicolumn{1}{l}{$\gamma$} & Coefficient of Gaussian kernel of LARR & $0.5$ \\
    \multicolumn{1}{l}{$l$} & Gaussian kernel size of LARR & $9$ \\
    \multicolumn{1}{l}{$\alpha$} & Update momentum of LARR & $0.9$ \\
    \multicolumn{1}{l}{$n_{bs}$} & Batch size & $32$ \\
    \multicolumn{1}{l}{$n_{epochs}$} & Maximum training epochs & $100$ \\
    \multicolumn{1}{l}{$\eta_{0}$} & Initial learning rate & $1 \times 10^{-4}$ \\
    \multicolumn{1}{l}{$p_{dropout}$} & Dropout rate & $0.01$ \\
    \bottomrule
    \end{tabular}%
\end{threeparttable}
}
\label{tab:hyperparameters_setting}%
\end{table}%

\subsubsection{Evaluation Metrics}

We employ four metrics, namely Mean Absolute Error (MAE), Mean Square Error (MSE), Root Mean Square Error (RMSE), and Coefficient of Determination ($R^{2}$), to evaluate the model's generalization capability for OOD degradation patterns. Specifically, lower values of MAE, MSE, and RMSE signify improved generalization performance, whereas higher values of $R^{2}$ indicate stronger generalization ability.

\begin{figure*}[h]
\captionsetup[subfloat]{font=scriptsize}
\centering
% --- Row 1 ---
\subfloat[iTransformer]{\includegraphics[width=1.35in, height=1.137in]{fig4_iTransformer.eps}%
    \label{fig:rul_itransformer}} % Optional label for subfigure
% \hfill % Distribute space horizontally
\subfloat[PatchTST]{\includegraphics[width=1.35in, height=1.137in]{fig4_PatchTST.eps}%
    \label{fig:rul_patchtst}}
% \hfill
\subfloat[TSMixer]{\includegraphics[width=1.35in, height=1.137in]{fig4_TSMixer.eps}%
    \label{fig:rul_tsmixer}}
% \hfill
\subfloat[MambaLithium]{\includegraphics[width=1.35in, height=1.137in]{fig4_MambaLithium.eps}%
    \label{fig:rul_mamba}}
% \hfill
\subfloat[PINN]{\includegraphics[width=1.35in, height=1.137in]{fig4_PINN.eps}%
    \label{fig:rul_pinn}}

% --- Row 2 ---
\subfloat[iTransformer$^\ast$]{\includegraphics[width=1.35in, height=1.137in]{fig4_iTransformer+.eps}%
    \label{fig:rul_itransformer_plus}}
% \hfill
\subfloat[PatchTST$^\ast$]{\includegraphics[width=1.35in, height=1.137in]{fig4_PatchTST+.eps}%
    \label{fig:rul_patchtst_plus}}
% \hfill
\subfloat[TSMixer$^\ast$]{\includegraphics[width=1.35in, height=1.137in]{fig4_TSMixer+.eps}%
    \label{fig:rul_tsmixer_plus}}
% \hfill
\subfloat[CRNN]{\includegraphics[width=1.35in, height=1.137in]{fig4_CRNN.eps}%
    \label{fig:rul_crnn}}
% \hfill
\subfloat[PINN$^\ast$]{\includegraphics[width=1.35in, height=1.137in]{fig4_PINN+.eps}%
    \label{fig:rul_pinn_plus}}

% --- Row 3 ---
\subfloat[DeNN-bagging]{\includegraphics[width=1.35in, height=1.137in]{fig4_bagging.eps}%
    \label{fig:rul_bagging}}
% \hfill
\subfloat[DeNN-fusion]{\includegraphics[width=1.35in, height=1.137in]{fig4_fusion.eps}%
    \label{fig:rul_fusion}}
% \hfill
\subfloat[DeNN-voting]{\includegraphics[width=1.35in, height=1.137in]{fig4_voting.eps}%
    \label{fig:rul_voting}}
% \hfill
\subfloat[ConvMMD]{\includegraphics[width=1.35in, height=1.137in]{fig4_ConvMMD.eps}%
    \label{fig:rul_convmmd}}
% \hfill
\subfloat[DRNet]{\includegraphics[width=1.35in, height=1.137in]{fig4_DRNet.eps}%
    \label{fig:rul_drnet}}

% --- Row 4 ---
\subfloat[DDAN]{\includegraphics[width=1.35in, height=1.137in]{fig4_DDAN.eps}%
    \label{fig:rul_ddan}}
\subfloat[DKTL]{\includegraphics[width=1.35in, height=1.137in]{fig4_DKTL.eps}%
    \label{fig:rul_dktl}}
\subfloat[LSTM-DA]{\includegraphics[width=1.35in, height=1.137in]{fig4_LSTM-DA.eps}%
    \label{fig:rul_lstmda}}
\subfloat[MAGNet]{\includegraphics[width=1.35in, height=1.137in]{fig4_MAGNet.eps}%
    \label{fig:rul_magnet}}
% Consider changing 'OURS' to your actual model name, e.g., 'RODTNet'
\subfloat[RODTNet]{\includegraphics[width=1.35in, height=1.137in]{fig4_RODTNet.eps}%
    \label{fig:rul_ours}}

\caption{Comparison of RUL prediction performance between the proposed RODTNet and comparison models on OOD generalization Task 1. In each panel, Predicted RUL (y-axis) is plotted against True RUL (x-axis). The dashed line serves as the ideal reference, and points represent individual predictions for the respective method. The distribution of prediction errors is visualized by the inset histogram located in the bottom right corner.}
\label{fig:rul} % Overall figure label
\end{figure*}

\begin{figure*}[b!]
\captionsetup[subfloat]{font=scriptsize}
\centering
\subfloat[Generalization Performance Comparison \label{fig:soh}]{\adjustbox{valign=t}{\includegraphics[height=5.5cm, keepaspectratio]{fig5_soh_radar_chart.eps}}}
\hspace{0.05\textwidth}
\subfloat[Generalization Efficiency Comparison \label{fig:comp}]{\adjustbox{valign=t}{\includegraphics[height=5.5cm, keepaspectratio]{fig5_footprint.eps}}}
\caption{Comparative analysis of the proposed RODTNet-LT model against baseline methods for SOH estimation.
(a) Radar chart comparing the performance of RODTNet-LT and baseline models across 11 different SOH estimation generalization tasks.
(b) Comparison of model performance (MSE), inference speed (s), and parameter count on the OOD SOH estimation generalization task.}
\label{fig:soh_efficiency}
\end{figure*}

% RUL对比实验
\begin{table*}[htbp]
    \footnotesize
    \centering
    \caption{Comparison experimental results for the OOD generalization tasks of RUL prediction on the MATR, TRI, TJU, and XJTU datasets.}
    % \resizebox{\linewidth}{!}{
    \scalebox{0.8}{
    \small
    \begin{adjustbox}{center}{
    \begin{threeparttable}
        \begin{tabular}{lccccccccccc}
        \toprule
        Task  & \multicolumn{2}{c}{Task1} & \multicolumn{2}{c}{Task2} & \multicolumn{2}{c}{Task3} & \multicolumn{2}{c}{Task4} & \multicolumn{2}{c}{Task5} & Task6 \\
    \cmidrule{2-12}    Metric & MAE   & MSE   & MAE   & MSE   & MAE   & MSE   & MAE   & MSE   & MAE   & MSE   & MAE \\
        \midrule
        iTransformer \cite{liu2024itransformer} & 99.78  & 15290.60  & 78.56  & 12430.67  & 153.83  & 34106.78  & 94.88  & 15629.50  & 93.21  & 16692.93  & 46.35  \\
        iTransformer$^\ast$\tnote{a} \cite{liu2024itransformer} & 12.16  & 290.75  & 44.35  & 5100.32  & 40.20  & 4450.44  & 29.83  & 3284.53  & 40.36  & 5311.12  & 9.80  \\
        PatchTST \cite{nie2023a} & 30.86  & 2147.96  & 141.59  & 40006.27  & 89.28  & 13432.70  & 149.46  & 36100.57  & 75.57  & 12738.93  & 29.58  \\
        PatchTST$^\ast$ \cite{nie2023a} & 12.82  & 393.49  & 42.24  & 4470.77  & 51.99  & 6520.95  & 38.72  & 3993.89  & 49.47  & 6372.25  & 30.24  \\
        TSMixer \cite{chen2023tsmixer} & 261.39  & 95673.01  & 95.70  & 17312.12  & 159.98  & 42316.82  & 125.36  & 27480.27  & 92.97  & 17272.87  & 115.18  \\
        TSMixer$^\ast$ \cite{chen2023tsmixer} & 16.77  & 584.11  & 52.31  & 5986.63  & 80.59  & 12311.22  & 75.14  & 12132.82  & 66.54  & 9945.74  & 19.64  \\
        PINN \cite{wang2024physics}  & 27.07  & 1345.26  & 29.33  & 1948.74  & 46.42  & 5144.30  & 46.24  & 4323.02  & 48.57  & 3539.95  & 43.84  \\
        PINN$^\ast$ \cite{wang2024physics} & 6.04  & 77.55  & 8.38  & 213.91  & 12.02  & 266.98  & 12.27  & 408.44  & 9.74  & 195.91  & 9.74  \\
        DR-Net \cite{WANG2023108897} & 11.61  & 284.32  & 42.50  & 5231.99  & 49.33  & 5978.47  & 42.87  & 5410.01  & 45.92  & 8061.80  & 8.91  \\
        CRNN \cite{D2EE01676A}  & 13.66  & 378.74  & 18.16  & 708.65  & 14.40  & 535.36  & 13.86  & 579.07  & 12.73  & 365.98  & 8.34  \\
        MAGNet \cite{TAN2024103725} & 13.51  & 347.84  & 15.08  & 492.65  & 16.03  & \underline{455.34}  & \underline{14.75}  & \underline{394.57}  & 13.37  & \underline{313.13}  & 8.98  \\
        DeNN-bagging \cite{10102497} & 23.40  & 1014.55  & 46.38  & 4144.07  & 68.88  & 8058.54  & 93.65  & 14047.69  & 57.23  & 5960.93  & 16.73  \\
        DeNN-fusion \cite{10102497} & 24.34  & 1010.97  & 49.15  & 4672.25  & 74.74  & 9549.55  & 95.93  & 16148.70  & 76.95  & 10290.79  & 16.80  \\
        DeNN-voting \cite{10102497} & 22.90  & 939.58  & 45.83  & 4074.09  & 71.58  & 8787.72  & 93.89  & 13623.54  & 57.07  & 6028.56  & 16.38  \\
        MambaLithium \cite{shi2024mambalithiumselectivestatespace} & 13.19  & 368.84  & 20.32  & 818.35  & 19.99  & 923.46  & 17.47  & 672.48  & 13.55  & 385.23  & 7.28  \\
        LSTM-DA \cite{HAN2022230823} & 20.36  & 713.70  & 30.14  & 1484.85  & 30.15  & 1701.16  & 31.45  & 1644.90  & 20.28  & 769.22  & 18.70  \\
        ConvMMD \cite{9788003} & 23.95  & 989.20  & 36.96  & 2619.11  & 46.70  & 4542.61  & 57.68  & 6594.72  & 28.70  & 1960.54  & 17.03  \\
        DKTL \cite{10029904}  & 23.77  & 1005.67  & 38.93  & 3220.22  & 51.40  & 5403.52  & 62.88  & 7711.63  & 30.94  & 2252.70  & 15.71  \\
        DDAN \cite{9560040}  & 19.35  & 622.24  & 26.90  & 1206.61  & 25.29  & 1267.76  & 26.25  & 1252.40  & 17.79  & 541.77  & 14.07  \\
        RODTNet$^{\ast\ast}$\tnote{b} {(ours)} & \underline{2.13}  & \underline{20.29}  & \underline{12.01}  & \underline{380.16}  & \underline{10.49}  & 577.50  & 25.63  & 3137.51  & \underline{12.83}  & 481.95  & \underline{4.06}  \\
        RODTNet (ours) & \textbf{0.53} & \textbf{1.36} & \textbf{2.38} & \textbf{25.62} & \textbf{1.94} & \textbf{31.62} & \textbf{1.95} & \textbf{97.98} & \textbf{3.90} & \textbf{85.62} & \textbf{0.55} \\
    \midrule
        \multicolumn{12}{c}{} \\
    \midrule
            & Task6 & \multicolumn{2}{c}{Task7} & \multicolumn{2}{c}{Task8} & \multicolumn{2}{c}{Task9} & \multicolumn{2}{c}{Task10} & \multicolumn{2}{c}{Task11} \\
    \cmidrule{2-12}          & MSE   & MAE   & MSE   & MAE   & MSE   & MAE   & MSE   & MAE   & MSE   & MAE   & MSE \\
        \midrule
        iTransformer \cite{liu2024itransformer} & 3348.46  & 79.31  & 11828.38  & 285.14  & 134288.68  & 97.72  & 12857.57  & 115.28  & 16008.10  & 86.33  & 9092.28  \\
        iTransformer$^\ast$ \cite{liu2024itransformer} & 451.88  & 19.96  & 690.05  & 25.39  & 1232.29  & 12.37  & 255.54  & 5.93  & 50.74  & 21.45  & 746.27  \\
        PatchTST \cite{nie2023a} & 1881.96  & 156.53  & 55065.22  & 171.01  & 70146.16  & 80.18  & 8273.63  & 78.84  & 8451.20  & 108.02  & 16314.58  \\
        PatchTST$^\ast$ \cite{nie2023a} & 2642.04  & 21.09  & 803.19  & 107.18  & 20704.56  & 13.70  & 298.71  & 7.34  & 81.72  & 25.52  & 1018.20  \\
        TSMixer \cite{chen2023tsmixer} & 21255.96  & 102.14  & 18803.24  & 363.95  & 183759.79  & 343.38  & 119826.16  & 97.57  & 12165.74  & 356.55  & 134948.97  \\
        TSMixer$^\ast$ \cite{chen2023tsmixer} & 844.83  & 20.28  & 700.89  & 45.71  & 3708.16  & 17.68  & 530.08  & 7.90  & 105.46  & 35.20  & 2180.88  \\
        PINN \cite{wang2024physics}  & 4496.71  & 81.96  & 18077.43  & 126.97  & 25586.95  & 53.12  & 5147.91  & 55.81  & 4744.68  & 111.30  & 22690.40  \\
        PINN$^\ast$ \cite{wang2024physics} & 195.91  & 16.70  & 610.41  & 24.96  & 1203.36  & 10.54  & 251.13  & 9.83  & 261.32  & 9.83  & 177.17  \\
        DR-Net \cite{WANG2023108897} & 180.23  & 15.77  & 482.02  & 30.98  & 1763.16  & 17.30  & 532.24  & 17.40  & 669.82  & 24.64  & 1306.43  \\
        CRNN \cite{D2EE01676A}  & 138.99  & 13.79  & 365.50  & 39.02  & 3541.77  & 29.13  & 1250.78  & 26.57  & 1346.65  & 35.58  & 3191.76  \\
        MAGNet \cite{TAN2024103725} & 146.94  & 15.79  & 428.60  & 44.58  & 4499.43  & 27.53  & 1183.21  & 24.84  & 1321.13  & 38.51  & 2961.67  \\
        DeNN-bagging \cite{10102497} & 607.03  & 28.97  & 1520.12  & 85.31  & 15854.20  & 40.39  & 2688.63  & 36.51  & 2682.08  & 61.02  & 7199.04  \\
        DeNN-fusion \cite{10102497} & 625.23  & 31.39  & 1769.45  & 92.75  & 17801.56  & 41.68  & 2863.98  & 36.61  & 2652.09  & 62.20  & 7491.24  \\
        DeNN-voting \cite{10102497} & 586.70  & 30.15  & 1646.72  & 86.36  & 16175.97  & 40.83  & 2692.80  & 36.85  & 2739.17  & 61.56  & 7397.65  \\
        MambaLithium \cite{shi2024mambalithiumselectivestatespace} & 105.15  & 21.18  & 752.52  & 71.95  & 12473.49  & 45.07  & 3804.16  & 28.68  & 1469.89  & 41.50  & 3260.10  \\
        LSTM-DA \cite{HAN2022230823} & 639.52  & 32.55  & 1741.70  & 121.14  & 24930.62  & 31.85  & 1526.07  & 24.87  & 1151.76  & 44.07  & 3391.50  \\
        ConvMMD \cite{9788003} & 544.12  & 37.69  & 2722.93  & 108.81  & 21251.08  & 20.64  & 668.67  & 28.49  & 1531.89  & 50.73  & 5750.57  \\
        DKTL \cite{10029904}  & 465.40  & 36.85  & 2282.25  & 104.67  & 19313.69  & 28.27  & 1281.96  & 27.02  & 1456.52  & 54.55  & 6445.32  \\
        DDAN \cite{9560040}  & 357.74  & 29.73  & 1478.19  & 117.50  & 24882.38  & 34.05  & 1859.57  & 24.13  & 1211.62  & 42.90  & 3599.68  \\
        RODTNet$^{\ast\ast}$ {(ours)} & \underline{36.10}  & \textbf{0.40} & \textbf{0.30} & \underline{1.51}  & \underline{6.89}  & \underline{7.62}  & \underline{127.05}  & \underline{3.86}  & \textbf{36.87} & \textbf{8.39} & \underline{419.11}  \\
        RODTNet (ours) & \textbf{0.68} & \underline{0.78}  & \underline{1.18}  & \textbf{0.32} & \textbf{0.29} & \textbf{4.34} & \textbf{46.06} & \textbf{3.50} & \underline{70.76}  & \underline{14.03}  & \textbf{395.27} \\
        \bottomrule
        \end{tabular}%}
    \begin{tablenotes}
    \item[a] $\ast$ indicates finetuning on the target domain.
    \item[b] $\ast\ast$ indicates training directly on the target domain.
    \item[c] \textbf{Bold}/\underline{underline} indicates the \textbf{best}/\underline{second} best.
    \item[d] Models utilizing pretraining-finetuning are excluded from comparison for fairness.
    \item[e] Notation and experimental settings described here also apply to Table~\ref{tab:soh_results}.
\end{tablenotes}
\end{threeparttable}}
\end{adjustbox}
}
\label{tab:rul_results}%
\end{table*}%

\begin{figure*}[htbp]
\vspace{-1em}
\centering
\subfloat{
    \includegraphics[width=1in]{fig6_rul_manifold_ours_tri_3d.eps}
    \includegraphics[width=1in]{fig6_rul_manifold_ours_tongji_3d.eps}
    \includegraphics[width=1in]{fig6_rul_manifold_ours_hust_3d.eps}
    \hfill
    \hspace{15pt} % 左侧间距加大
    \smash{\begin{tikzpicture}
        \draw[gray!40, densely dashed, line width=0.6pt] (0,0.1in) -- (0,0.9in); % 缩短高度并上下留空
    \end{tikzpicture}}
    \hspace{15pt} % 右侧间距加大
    \hfill
    \includegraphics[width=1in]{fig6_rul_manifold_others_tri_3d.eps}
    \includegraphics[width=1in]{fig6_rul_manifold_others_tongji_3d.eps}
    \includegraphics[width=1in]{fig6_rul_manifold_others_hust_3d.eps}
}
\hfil
\subfloat{
    \includegraphics[width=1in]{fig6_rul_manifold_ours_tri_2d.eps}
    \includegraphics[width=1in]{fig6_rul_manifold_ours_tongji_2d.eps}
    \includegraphics[width=1in]{fig6_rul_manifold_ours_hust_2d.eps}
    \hfill
    \hspace{15pt}
    \smash{\begin{tikzpicture}
        \draw[gray!40, densely dashed, line width=0.6pt] (0,0.1in) -- (0,0.9in);
    \end{tikzpicture}}
    \hspace{15pt}
    \hfill
    \includegraphics[width=1in]{fig6_rul_manifold_others_tri_2d.eps}
    \includegraphics[width=1in]{fig6_rul_manifold_others_tongji_2d.eps}
    \includegraphics[width=1in]{fig6_rul_manifold_others_hust_2d.eps}
}
\hfil
\subfloat{
    \includegraphics[width=1in]{fig6_rul_manifold_ours_tri_1d.eps}
    \includegraphics[width=1in]{fig6_rul_manifold_ours_tongji_1d.eps}
    \includegraphics[width=1in]{fig6_rul_manifold_ours_hust_1d.eps}
    \hfill
    \hspace{15pt}
    \smash{\begin{tikzpicture}
        \draw[gray!40, densely dashed, line width=0.6pt] (0,0.1in) -- (0,0.9in);
    \end{tikzpicture}}
    \hspace{15pt}
    \hfill
    \includegraphics[width=1in]{fig6_rul_manifold_others_tri_1d.eps}
    \includegraphics[width=1in]{fig6_rul_manifold_others_tongji_1d.eps}
    \includegraphics[width=1in]{fig6_rul_manifold_others_hust_1d.eps}
}
\caption{Comparison of feature distributions obtained from RODTNet (left) and TSMixer (right). Features are visualized via PCA dimensionality reduction to 3D (top row), 2D (middle row), and 1D (bottom row) and points are colored based on ground truth labels.}
\label{fig:Performance_Comparison_2}
\end{figure*}

% % SOH对比实验
% \begin{table*}[htbp]
%     \centering
%     \caption{Comparison experimental results for the OOD generalization tasks of SOH estimation on the MATR, TRI, TJU, and XJTU datasets.}
%     % \resizebox{\linewidth}{!}{
%     \scalebox{0.8}{
%     \small
%     \begin{threeparttable}
%     \begin{tabular}{lccccccccccc}
%     \toprule
%     Task  & \multicolumn{2}{c}{Task1} & \multicolumn{2}{c}{Task2} & \multicolumn{2}{c}{Task3} & \multicolumn{2}{c}{Task4} & \multicolumn{2}{c}{Task5} & \multicolumn{1}{l}{Task6} \\
% \cmidrule{2-12}    Metric & MAE   & MSE   & MAE   & MSE   & MAE   & MSE   & MAE   & MSE   & MAE   & MSE   & MAE \\
%     \midrule
%     iTransformer \cite{liu2024itransformer} & 5.9093  & 0.5257  & 1.1315  & 0.0350  & 1.3424  & 0.0317  & 1.0941  & 0.0218  & 2.7314  & 0.1079  & 21.3182  \\
%     iTransformer$^\ast$ \cite{liu2024itransformer} & 0.7803  & 0.0134  & 0.7618  & 0.0167  & 0.4686  & 0.0067  & 0.4534  & 0.0051  & 0.7615  & 0.0183  & 2.1108  \\
%     PatchTST \cite{nie2023a} & \underline{1.1228}  & 0.0311  & 4.1024  & \underline{0.3776}  & 0.7981  & 0.0145  & 4.2731  & 0.4299  & 2.8899  & 0.1764  & 9.5076  \\
%     PatchTST$^\ast$ \cite{nie2023a} & 0.6310  & 0.0098  & 0.6937  & 0.0136  & 0.4426  & 0.0071  & 0.4094  & 0.0046  & 0.8018  & 0.0183  & 3.6897  \\
%     TSMixer \cite{chen2023tsmixer} & 16.2029  & 3.4749  & 2.5634  & 0.1033  & 5.1793  & 0.3115  & 2.9573  & 0.0980  & 3.1299  & 0.1231  & 35.9175  \\
%     TSMixer$^\ast$ \cite{chen2023tsmixer} & 0.8600  & 0.0172  & 0.7558  & 0.0160  & 0.5797  & 0.0082  & 0.4818  & 0.0046  & 0.8081  & 0.0184  & 2.9360  \\
%     PINN \cite{wang2024physics}  & 1.3596  & 0.0403  & 2.2216  & 0.0844  & 1.4187  & 0.0511  & 1.1438  & 0.0227  & 1.7547  & 0.0556  & 20.7705  \\
%     PINN$^\ast$ \cite{wang2024physics} & 0.5748  & 0.0061  & 1.3991  & 0.0361  & 1.3524  & 0.0269  & 0.7226  & 0.0107  & 1.0708  & 0.0207  & 8.9555  \\
%     DR-Net \cite{WANG2023108897} & 4.2229  & 0.3324  & 2.6507  & 0.1112  & 3.4879  & 0.1749  & 3.6984  & 0.2049  & 3.4467  & 0.1925  & 16.0808  \\
%     CRNN \cite{D2EE01676A}  & 0.8260  & 0.0168  & 0.5589  & 0.0058  & 0.5295  & 0.0060  & 0.5508  & 0.0082  & 0.8301  & 0.0184  & 7.6085  \\
%     MAGNet \cite{TAN2024103725} & 1.2410  & \underline{0.0244}  & 0.9305  & 0.0146  & 1.1954  & 0.0233  & 0.9216  & 0.0196  & 1.8908  & 0.0526  & 6.9453  \\
%     DeNN-bagging \cite{10102497} & 1.7015  & 0.0658  & 2.1448  & 0.1194  & 2.1157  & 0.1141  & 2.1386  & 0.1330  & 2.8398  & 0.1719  & 10.3293  \\
%     DeNN-fusion \cite{10102497} & 1.5901  & 0.0542  & 2.2918  & 0.1403  & 2.1127  & 0.1162  & 2.3360  & 0.1544  & 4.0800  & 0.2782  & 10.3575  \\
%     DeNN-voting \cite{10102497} & 1.6199  & 0.0577  & 2.1824  & 0.1217  & 2.0499  & 0.1055  & 2.1529  & 0.1109  & 2.8789  & 0.1806  & 10.1680  \\
%     MambaLithium \cite{shi2024mambalithiumselectivestatespace} & 1.2177  & 0.0338  & \underline{0.4801}  & 0.0038  & \underline{0.4002}  & \underline{0.0031}  & \underline{0.4162}  & \underline{0.0035}  & \underline{0.9160}  & \underline{0.0258}  & \underline{5.5906}  \\
%     LSTM-DA \cite{HAN2022230823} & 3.0477  & 0.1140  & 2.8469  & 0.0945  & 3.1725  & 0.1164  & 3.4704  & 0.1360  & 3.4921  & 0.1370  & 12.8066  \\
%     ConvMMD \cite{9788003} & 2.2734  & 0.0891  & 2.2941  & 0.0809  & 2.4317  & 0.0862  & 2.6367  & 0.0943  & 2.5047  & 0.0807  & 10.7025  \\
%     DKTL \cite{10029904}  & 2.2029  & 0.0889  & 2.4482  & 0.1134  & 2.6035  & 0.1153  & 2.8678  & 0.1272  & 2.7184  & 0.1135  & 10.5311  \\
%     DDAN \cite{9560040}  & 2.2708  & 0.0667  & 1.7973  & 0.0378  & 2.0361  & 0.0485  & 2.3470  & 0.0611  & 2.4232  & 0.0648  & 12.1931  \\
%     RODTNet-LT$^{\ast\ast}$ {(ours)} & 1.1934  & 0.0484  & 0.9942  & 0.0367  & 0.9038  & 0.0330  & 2.4697  & 0.3593  & 3.1390  & 0.4989  & 9.6180  \\
%     RODTNet-LT (ours) & \textbf{1.0384} & \textbf{0.0208} & \textbf{0.4476} & \textbf{0.0030} & \textbf{0.3630} & \textbf{0.0020} & \textbf{0.2885} & \textbf{0.0014} & \textbf{0.4435} & \textbf{0.0032} & \textbf{4.0520} \\
%     \midrule
%     \multicolumn{12}{c}{} \\
%     \midrule
%             & \multicolumn{1}{l}{Task6} & \multicolumn{2}{c}{Task7} & \multicolumn{2}{c}{Task8} & \multicolumn{2}{c}{Task9} & \multicolumn{2}{c}{Task10} & \multicolumn{2}{c}{Task11} \\
% \cmidrule{2-12}          & MSE   & MAE   & MSE   & MAE   & MSE   & MAE   & MSE   & MAE   & MSE   & MAE   & MSE \\
%     \midrule
%     iTransformer \cite{liu2024itransformer} & 8.5517  & 10.6146  & 1.6506  & 9.1256  & 1.5742  & 21.2702  & 8.0010  & 40.3827  & 20.7572  & 19.2927  & 4.6340  \\
%     iTransformer$^\ast$ \cite{liu2024itransformer} & 0.1204  & 8.8424  & 1.4328  & 1.3233  & 0.0277  & 2.9985  & 0.1737  & 2.8956  & 0.1635  & 2.8476  & 0.1570  \\
%     PatchTST \cite{nie2023a} & 1.8668  & 16.9823  & 5.4530  & 9.2243  & 1.6822  & 14.7543  & 4.8261  & 48.1532  & 28.1815  & 15.8620  & 4.5693  \\
%     PatchTST$^\ast$ \cite{nie2023a} & 0.5973  & 10.1355  & 1.8044  & 6.8517  & 0.9185  & 3.3279  & 0.1986  & 4.6602  & 0.3933  & 5.0112  & 0.4733  \\
%     TSMixer \cite{chen2023tsmixer} & 15.7721  & 8.1766  & 1.0106  & 14.3296  & 2.9821  & 51.0444  & 26.6658  & 44.3269  & 22.9838  & 58.9715  & 35.8934  \\
%     TSMixer$^\ast$ \cite{chen2023tsmixer} & 0.1595  & 9.3981  & 1.4881  & 1.6831  & 0.0443  & 2.7117  & 0.1365  & 4.0151  & 0.2957  & 3.7528  & 0.2954  \\
%     PINN \cite{wang2024physics}  & 6.2435  & 25.2663  & 9.1038  & 5.9623  & 0.7837  & 7.9811  & 1.9312  & 22.6493  & 7.5184  & 13.4784  & 3.3358  \\
%     PINN$^\ast$ \cite{wang2024physics} & 2.2297  & 8.1476  & 1.3071  & 1.3870  & 0.0357  & 3.9588  & 0.4028  & 5.6329  & 0.9141  & 3.6541  & 0.3709  \\
%     DR-Net \cite{WANG2023108897} & 3.9837  & 19.8591  & 5.6675  & 40.7034  & 21.2569  & 13.3021  & 3.0780  & 10.6609  & 2.0293  & 10.4743  & 1.7780  \\
%     CRNN \cite{D2EE01676A}  & 1.6485  & 4.5716  & 0.4770  & 3.3587  & 0.3612  & 5.8487  & 0.8072  & 11.0377  & 3.0051  & 6.9622  & 1.2084  \\
%     MAGNet \cite{TAN2024103725} & 1.2605  & \underline{3.6287}  & \textbf{0.2369} & 3.2225  & 0.2692  & \underline{3.4654}  & \underline{0.2234}  & 7.5763  & 1.7185  & \underline{4.1481}  & \underline{0.4131}  \\
%     DeNN-bagging \cite{10102497} & 1.9811  & 7.1896  & 0.8327  & 5.6580  & 0.7403  & 7.6599  & 1.1013  & 12.8670  & 3.4043  & 8.8430  & 1.6496  \\
%     DeNN-fusion \cite{10102497} & 1.9799  & 8.5083  & 1.0813  & 5.9468  & 0.7503  & 8.3443  & 1.3646  & 12.5343  & 3.2910  & 8.9251  & 1.7611  \\
%     DeNN-voting \cite{10102497} & 1.9101  & 7.4300  & 0.8838  & 5.5938  & 0.7390  & 7.7196  & 1.1244  & 12.7772  & 3.4203  & 8.9074  & 1.6948  \\
%     MambaLithium \cite{shi2024mambalithiumselectivestatespace} & \underline{0.7132}  & 6.7626  & 0.8283  & 4.8486  & 0.6418  & 7.0512  & 1.2408  & 8.7161  & 2.1239  & 5.8056  & 0.8398  \\
%     LSTM-DA \cite{HAN2022230823} & 2.8925  & 8.3498  & 1.0564  & 7.9172  & 1.3040  & 6.0889  & 0.6520  & 10.2645  & 2.3890  & 7.8039  & 1.1160  \\
%     ConvMMD \cite{9788003} & 2.0985  & 7.8693  & 0.9537  & 6.7908  & 0.9026  & 5.9231  & 0.6199  & 11.1645  & 2.6587  & 7.6832  & 1.2462  \\
%     DKTL \cite{10029904}  & 2.1718  & 7.5353  & 0.8587  & 6.7470  & 0.9272  & 4.9570  & 0.5125  & 10.8147  & 2.5730  & 6.9723  & 1.1317  \\
%     DDAN \cite{9560040}  & 2.7491  & 8.1399  & 1.1031  & 7.3573  & 1.1792  & 5.9360  & 0.8296  & 9.0353  & 2.1172  & 5.9290  & 0.8442  \\
%     RODTNet-LT$^{\ast\ast}$ {(ours)} & 2.0079  & 4.2455  & 0.4090  & \underline{1.7140}  & \underline{0.0635}  & \textbf{3.0794} & \textbf{0.2118} & \textbf{3.3261} & \textbf{0.2229} & \textbf{3.6471} & 0.8305  \\
%     RODTNet-LT (ours) & \textbf{0.3946} & \textbf{3.5394} & \underline{0.3020}  & \textbf{1.0781} & \textbf{0.0195} & 4.5385  & 0.3831  & \underline{5.0415}  & \underline{0.3875}  & 4.3266  & \textbf{0.3424} \\
%     \bottomrule
%     \end{tabular}%}
%     %   \begin{tablenotes}
%     %     \item[a] $\ast$ indicates finetuning on the target domain.
%     %     \item[b] $\ast\ast$ indicates training directly on the target domain.
%     %     \item[c] \textbf{Bold}/\underline{underline} indicates the \textbf{best}/\underline{second} best
%     %     \item[d] Models utilizing pretraining-finetuning are excluded from comparison for fairness.
%     %     \item[e] Notation and experimental settings described here also apply to Table~\ref{tab:soh_results}.
%     % \end{tablenotes}
% \end{threeparttable}
% }
%     \label{tab:soh_results}%
% \end{table*}%

\subsection{Performance Comparison on OOD Generalization}\label{subsec:comparative_experiments}

To evaluate the OOD generalization capability of our proposed RODTNet and RODTNet-LT, we conduct comprehensive experiments across 11 tasks using MAE and MSE as evaluation metrics. For RUL prediction tasks, Table~\ref{tab:rul_results} demonstrates that RODTNet achieves state-of-the-art performance in 9 tasks and ranks second in the remaining 2 tasks, surpassed only by its target-domain-adapted variant.
% 
% In SOH estimation tasks, RODTNet-LT attains optimal performance in 9 scenarios as shown in Table~\ref{tab:soh_results}, while exhibiting marginal performance gaps against its target-domain counterpart and MAGNet in 2 cases.
The complete SOH estimation results are provided in \textbf{Table S1 of the Supplementary Material}. Notably, RODTNet-LT achieves the best performance in 9 out of 11 scenarios, showing only a marginal gap compared to its target-domain counterpart and MAGNet in 2 cases.
% 
Quantitatively, in OOD RUL prediction tasks, RODTNet reduces prediction errors by 40.76\% MAE and 28.65\% MSE compared to suboptimal baselines. Similarly, in OOD SOH estimation tasks, RODTNet-LT achieves 5.03\% MAE and 22.22\% MSE improvements over existing benchmarks. These average percentage improvements are calculated by comparing our method's performance against the second-best baseline when it is optimal for a given task, or against the best-performing baseline otherwise (with negative values denoting underperformance relative to the best), and then averaging these relative differences across all 11 generalization tasks. These enhancements validate the proposed method's robustness against domain shifts induced by heterogeneous degradation mechanisms.
% 
In contrast, existing baselines demonstrate limited capability in extracting domain-invariant degradation patterns, even under target domain finetuning/adaptation with partial data access. Notably, both RODTNet and RODTNet-LT outperform their target-domain-trained variants, which can be attributed to the enhanced pattern comprehensiveness achieved through multi-source domain integration.
% 
Visual evidence further corroborates these findings.
Fig.~\ref{fig:rul} assesses RODTNet's generalization against baselines using scatter plots (actual vs. predicted life) and error distributions histograms. Each subplot also presents the $R^{2}$ score in the top-left corner, further quantifying the goodness of fit. RODTNet's predictions cluster densely around the ideal line with a significantly narrower error distribution, demonstrating its effectiveness in capturing OOD degradation patterns. Furthermore, as shown by the radar charts in Fig.~\ref{fig:soh_efficiency}(a) for SOH estimation, the performance polygon for RODTNet-LT consistently extends closer to the outer edges of the chart across multiple metrics, reflecting its superior and more comprehensive generalization performance.

To investigate the topology of degradation representations, we apply PCA for dimensionality reduction to three, two, and one dimensions to visualize feature maps from the RODTNet encoder, coloring points by labels. \rev{As depicted in Fig.~\ref{fig:Performance_Comparison_2}, the manifold obtained with the MCRDR strategy on the left exhibits a continuous distribution, well-preserved across dimensions, facilitating high-precision generalization. Such topological continuity ensures smooth outputs, benefiting downstream BMS control stability.} Conversely, the manifold without MCRDR, shown on the right, displays discontinuities appearing as holes and suffers from significant feature overlap between degradation stages upon dimensionality reduction, hindering such generalization.

An evaluation of the trade-off between prognostic performance, inference speed, and model size for RODTNet-LT and baseline methods is presented, with results visualized in Fig.~\ref{fig:soh_efficiency}(b). The findings show that models achieving prognostic performance comparable to RODTNet-LT, such as MambaLithium and MAGNet, tend to be larger in size or slower in inference. Conversely, models with greater parameter or time efficiency, including PINN, ConvMMD, and TSMixer, exhibit inferior generalization. This underscores RODTNet-LT's effective balance, which stems from its design. On one hand, its heuristic signal reconstruction yields compact degradation features. On the other, its Mamba encoder offers linear complexity, and the uncertainty-driven feature manipulation pre-computes manifold boundary extrapolations and feature corrections. \rev{Thus, RODTNet-LT achieves robust generalization with a low computational footprint, verifying its real-time feasibility on automotive electronic control units.}

% ablation study 1
\begin{table*}[t!]
    \footnotesize
    \centering
    \caption{Ablation study results of RODTNet on the MATR, Tongji, HUST, and XJTU datasets.}
    \scalebox{0.8}
    {
        \small
        \begin{threeparttable}
        \begin{tabular}{lccccccccccc}
        \toprule
            & \multicolumn{3}{c}{Component} & \multicolumn{2}{c}{MATR Dataset} & \multicolumn{2}{c}{Tongji Dataset} & \multicolumn{2}{c}{HUST Dataset} & \multicolumn{2}{c}{XJTU Dataset} \\
    \cmidrule{2-12}          & \multicolumn{1}{l}{MC} & \multicolumn{1}{l}{Mamba} & \multicolumn{1}{l}{UDFS} & MAE   & MSE   & MAE   & MSE   & MAE   & MSE   & MAE   & MSE \\
        \midrule
        Baseline\tnote{a} &  &  &  & 111.2889  & 17628.7479  & 21.0490  & 926.8474  & 266.2715  & 116654.9141  & 96.2077  & 12405.8422  \\
        DGBL\tnote{b} &       &       & \checkmark & 23.3242  & 979.1668  & 22.1563  & 852.3195  & 65.1147  & 6879.1396  & 32.4785  & 1900.7478  \\
        LMCRDR\tnote{c} & \checkmark &       &       & 0.9910  & 1.8321  & 3.2627  & 21.1491  & 0.9775  & 2.8166  & 9.5687  & 163.8348  \\
        MCRDR\tnote{d} & \checkmark & \checkmark & & \underline{0.6214}  & 1.3798  & \underline{0.8529}  & \underline{1.4767}  & \underline{0.4390}  & \underline{0.4802}  & \underline{4.7762}  & \underline{49.0050}  \\
        LRODTNet\tnote{e} & \checkmark &  & \checkmark & 0.6734  & \textbf{0.7131} & 2.9669  & 16.7873  & 1.3954  & 3.9988  & 10.6988  & 234.1704  \\
        RODTNet & \checkmark & \checkmark & \checkmark & \textbf{0.5370} & \underline{1.3636}  & \textbf{0.7820} & \textbf{1.1815} & \textbf{0.3200} & \textbf{0.2935} & \textbf{4.3448} & \textbf{46.0641} \\
        \bottomrule
        \end{tabular}%
        \begin{tablenotes}
        \item[a] Baseline: Model with downsampled raw cycling signals and linear backbone.
        \item[b] DGBL: Baseline model enhanced with feature manipulation strategy.
        \item[c] LMCRDR: RODTNet variant with linear backbone, without feature manipulation.
        \item[d] MCRDR: RODTNet variant without the feature manipulation strategy.
        \item[e] LRODTNet: RODTNet variant with linear backbone.
    \end{tablenotes}
    \end{threeparttable}
    }
    \label{tab:ablation_1}%
    \end{table*}%
    
    % ablation study 2
    \begin{table}[htbp]
    \centering
    \caption{Ablation study results of RODTNet-LT on the MATR, Tongji, HUST, and XJTU datasets.}
    % \resizebox{\linewidth}{!}{
    \scalebox{0.7}{
    \small
    \begin{adjustbox}{center}
    {
        \begin{threeparttable}
        \begin{tabular}{lcccccccc}
        \toprule
        \multicolumn{1}{c}{Dataset} & \multicolumn{2}{c}{MATR Dataset} & \multicolumn{2}{c}{Tongji Dataset} & \multicolumn{2}{c}{HUST Dataset} & \multicolumn{2}{c}{XJTU Dataset} \\
    \cmidrule{2-9}    \multicolumn{1}{c}{Metric} & MAE   & MSE   & MAE   & MSE   & MAE   & MSE   & MAE   & MSE \\
        \midrule
        Baseline\tnote{a} & 5.1276  & 0.5883  & 9.6675  & 1.6127  & 9.4620  & 1.4738  & 20.8255  & 9.4026  \\
        RODTNet & \underline{1.1207}  & \underline{0.0257}  & \underline{7.3428}  & \underline{0.8164}  & \underline{1.5201}  & \underline{0.0448}  & \textbf{4.4705} & \underline{0.3998}  \\
        RODTNet-LT\tnote{b}  & \textbf{1.0384} & \textbf{0.0208} & \textbf{3.5394} & \textbf{0.3020} & \textbf{1.0781} & \textbf{0.0195} & \underline{4.5385}  & \textbf{0.3831} \\
        \bottomrule
        \end{tabular}%
        \begin{tablenotes}
        \item[a] Baseline: Model with downsampled raw cycling signals and linear backbone.
        \item[b] RODTNet-LT: RODTNet with LARR module.
    \end{tablenotes}
    \end{threeparttable}}
    \end{adjustbox}
    }
    \label{tab:ablation_study_2}%
    \end{table}%

\begin{figure*}[h]
    \vspace{-1em}
    \centering
    \includegraphics[width=7.16in]{Fig7_sensitivity_analysis_combined.eps}
    \caption{Hyperparameter sensitivity analysis of RODTNet's generalization performance evaluated by RMSE across the MATR, Tongji, HUST, and XJTU datasets, investigating the impact of varying batch sizes, dimensions of hidden state, the number of encoder layers, and USDF activation thresholds.}
    \label{fig:sensitivity}
    \end{figure*}

\subsection{Ablation Study}

We conduct ablation studies to evaluate the effectiveness of each components in RODTNet and RODTNet-LT. Table~\ref{tab:ablation_1} reports the MAE and MSE results for RUL prediction using RODTNet variants across 4 datasets. Correspondingly, Table~\ref{tab:ablation_study_2} presents these metrics for the SOH estimation task using both RODTNet and RODTNet-LT on the same datasets.

\subsubsection{Baseline}

We establish the baseline model using a fully linear backbone with downsampled raw cyclic signals as input.

\subsubsection{DGBL}

We obtain the DGBL model by adding UDFS strategy to the baseline. We observe an improvement in generalization performance, which proves the effectiveness of our proposed UDFS in expanding the distribution boundary and fitting OOD degradation patterns.

\subsubsection{LMCRDR}

We create the LMCRDR model by replacing the baseline's raw input with reconstructed input using the heuristic manifold contraction strategy. We observe an improvement in performance, demonstrating the good adaptability of our manifold contraction approach to the backbone and its superiority in generalizing degradation patterns.

\subsubsection{MCRDR}

We then develop the MCRDR model by replacing the linear backbone of LMCRDR with Mamba. We observed a further improvement in prognostic accuracy, indicating that Mamba effectively extracts domain-invariant features.

\subsubsection{LRODTNet}

We combine LMCRDR with the UDFS while retaining a linear backbone to form LRODTNet. We observe that the addition of UDFS improved performance over LMCRDR, thus proving that our propose feature synthesis strategy can extend the boundary of learned degradation representations possessing low intrinsic dimensions.

\subsubsection{RODTNet-LT}

For the task of SOH estimation involving long-tailed distributions, we obtained the RODTNet-LT variant by the adding LARR strategy. We observe that RODTNet-LT achieved lower generalization error compared to RODTNet, proving the effectiveness of LARR in removing head bias. \rev{This is vital for safety-critical monitoring during the high-risk battery end-of-life phase.}

\subsection{Hyperparameter Sensitivity Analysis}

We conduct a sensitivity analysis across 4 datasets to evaluate the model's performance sensitivity to variations in batch size, encoder hidden dimension, number of encoder layers, and the UDFS activation threshold, using target domain RMSE as the evaluation metric.

\subsubsection{Batch Size}

We investigate batch size's impact on performance, testing \{32, 64, 128, 256, 512\}. Fig.~\ref{fig:sensitivity}(a) demonstrates that a size of 32 typically yields optimal performance in general scenarios. However, on MATR tasks, larger batches decrease error, suggesting sufficient samples allow approximation of domain shift, whereas smaller batches better capture degradation patterns in other tasks.

\subsubsection{Encoder Hidden Dimension}

We assess how encoder hidden dimension affects generalization, testing \{16, 32, 64, 128, 256\}. Fig.~\ref{fig:sensitivity}(b) reveals improved performance on Tongji and HUST tasks with larger dimensions, reflecting scalability. Conversely, performance degrades on MATR and XJTU, suggesting potential overfitting from excessive model scale.

\subsubsection{Number of Encoder Layers}

We analyze the number of encoder layers' influence on model scale and generalization, testing \{1, 2, 3, 4, 5\}. Fig.~\ref{fig:sensitivity}(c) indicates effects generally consistent with hidden dimension, where increasing layers appear more susceptible to overfitting bottlenecks.

\subsubsection{UDFS Activation Threshold}

We evaluate sensitivity to the UDFS activation threshold, which affects the real/generated data ratio, testing \{0.0008, 0.004, 0.02, 0.1, 0.5\}. Fig.~\ref{fig:sensitivity}(d) shows generalization improves across all datasets as the threshold increases, highlighting the importance of real source domain data for modeling OOD patterns.

We observe RODTNet is generally insensitive to hyperparameters, except for the complex XJTU task. For consistent comparisons, we fix parameters across tasks using the frequently optimal batch size of 32, along with hidden dimension of 16 and encoder layer of 1 for computational efficiency. \rev{Notably, RODTNet maintains strong performance even with these efficiency-prioritized parameters, highlighting its design tailored for embedded vehicular systems.} We select a UDFS threshold of 0.5, effectively balancing minimal generalization error with reduced training overhead.  

\section{Conclusion}

\rev{In this paper, we have presented RODTNet, a novel method for robust OOD battery degradation modeling in vehicular BMS, along with its long-tail-optimized variant RODTNet-LT.}
% 
First, we have proposed a manifold contraction-based robust degradation representation extraction strategy. It achieves pattern continuity and compresses the feature manifold by combining heuristic signal reconstruction with latent dynamics generated via the Mamba mechanism.
% 
\rev{Second, we have developed an uncertainty-driven feature synthesis approach to address domain shifts arising from diverse automotive operating conditions.} This method correlates feature statistics uncertainty with domain variations using probabilistic modeling, steering the creation of domain-invariant features and diverse OOD examples.
% 
Finally, we have designed a long-tail-aware adaptive rebalancing refinement technique based on feature distribution smoothing. By exploiting feature-label correlations, it alleviates head-bias in the synthesized features, improving the handling of imbalanced trajectories.
% 
\rev{Comprehensive experimental results across 11 distinct OOD generalization tasks using 4 LIB aging datasets demonstrate the significant advantages of the proposed RODTNet and RODTNet-LT over 14 baselines and their 20 variants, advancing reliable vehicular BMS.}

% Despite advancements, real-world battery operation necessitates continuous adaptation as degradation dynamics evolve throughout the lifespan, posing demands beyond initial OOD robustness. Exploring temporal domain adaptation thus presents a promising future direction, aiming to equip models with dynamic online update capabilities for sustained prognostic performance over time.

\bibliographystyle{IEEEtran}
\bibliography{IEEEabrv, reference.bib}

% \newpage

% % If you have an EPS/PDF photo (graphicx package needed), extra braces are
% %  needed around the contents of the optional argument to biography to prevent
% %  the LaTeX parser from getting confused when it sees the complicated
% %  $\backslash${\tt{includegraphics}} command within an optional argument. (You can create
% %  your own custom macro containing the $\backslash${\tt{includegraphics}} command to make things
% %  simpler here.)
 
% % \vspace{11pt}

% % \bf{If you include a photo:}\vspace{-33pt}
% \begin{IEEEbiography}[{\includegraphics[width=1in,height=1.25in,clip,keepaspectratio]{fig1}}]{Michael Shell}
% Use $\backslash${\tt{begin\{IEEEbiography\}}} and then for the 1st argument use $\backslash${\tt{includegraphics}} to declare and link the author photo.
% Use the author name as the 3rd argument followed by the biography text.
% \end{IEEEbiography}

% \vfill

\end{document}
